\documentclass[11pt]{article}
\usepackage[T1]{fontenc}
\usepackage{lmodern}
\usepackage[margin=1in]{geometry}
\usepackage{amsmath,amssymb,amsthm,mathtools}
\usepackage{enumitem}
\usepackage{microtype}
\usepackage{xcolor}
\usepackage{hyperref}
\pdfinfoomitdate=1
\pdftrailerid{}
\pdfsuppressptexinfo=15
\hypersetup{colorlinks=true,linkcolor=black,citecolor=black,urlcolor=blue,
  pdftitle={Positivity of Serre's Intersection Multiplicity},pdfauthor={OpenAI}}
\setlength{\emergencystretch}{2em}
\numberwithin{equation}{section}
\newtheorem{theorem}{Theorem}[section]
\newtheorem{proposition}[theorem]{Proposition}
\newtheorem{lemma}[theorem]{Lemma}
\newtheorem{corollary}[theorem]{Corollary}
\theoremstyle{definition}
\newtheorem{definition}[theorem]{Definition}
\theoremstyle{remark}
\newtheorem{remark}[theorem]{Remark}
\newcommand{\Perf}{\operatorname{Perf}}
\newcommand{\Spec}{\operatorname{Spec}}
\newcommand{\Tor}{\operatorname{Tor}}
\newcommand{\rank}{\operatorname{rank}}
\newcommand{\length}{\operatorname{length}}
\newcommand{\Q}{\mathbb{Q}}
\newcommand{\Z}{\mathbb{Z}}
\newcommand{\m}{\mathfrak{m}}
\newcommand{\K}{K}
\newcommand{\cL}{\mathcal{L}}
\title{Positivity of Serre's Intersection Multiplicity}
\author{OpenAI}
\date{September 23, 2026}
\begin{document}
\maketitle
\begin{abstract}
We prove Serre's positivity conjecture, including ramified mixed
characteristic. If two nonzero finitely generated modules over a regular local
ring have tensor product of finite length and complementary dimensions, then
their intersection multiplicity is strictly positive.
\end{abstract}

\section{Introduction}\label{sec:introduction}

Let $(R,\m)$ be a commutative Noetherian regular local ring of dimension
$d$. For finite $R$-modules $M,N$ whose tensor product has finite length,
Serre's intersection multiplicity is the integer
\[
  \chi^R(M,N)=\sum_{i=0}^{d}(-1)^i
       \length_R\Tor_i^R(M,N).
\]
The support condition makes every term finite, and regularity makes the
higher Tor groups vanish. This Euler characteristic measures the local
intersection of the supports while accounting for higher Tor. Its
positivity is not a formal consequence of the nonzero degree-zero term:
the higher terms enter with alternating signs.

\begin{theorem}\label{thm:main}
Let $(R,\m)$ be a commutative Noetherian regular local ring of dimension
$d$, and let $M,N$ be nonzero finitely generated $R$-modules. If
$\length_R(M\otimes_R N)<\infty$ and
$\dim_R M+\dim_R N=d$, then
\[
  \chi^R(M,N)>0.
\]
\end{theorem}

Here regularity means that $\m$ can be generated by $d$ elements, and
module dimension means the Krull dimension of its support. Theorem
\ref{thm:main} proves Serre's positivity conjecture without a restriction
on characteristic or ramification. In particular, it does not require
either module to lift to an unramified regular local ring.\footnote{A
public post by Editan, dated December~9, 2025 and credited
to Editan, Copilot, and Gemini, also asserts the prime-quotient form of
Serre's positivity conjecture. We record this as a same-scope claim,
not as an established theorem; see
\href{https://note.com/editan/n/n35dec2c1ddbe}{the author's post}.}

\subsection{Background and significance}

Serre introduced intersection multiplicities through the Euler
characteristic of Tor. His reduction to the diagonal proves positivity
in equicharacteristic and in the unramified setting; the same argument
also treats formal power series rings over an arbitrary complete
discrete valuation ring \cite[Chapter~V]{Serre}. Thus ramification alone
does not separate the known cases from the general problem.
Roberts and, independently, Gillet--Soul\'e proved vanishing when the
tensor product has finite length and the sum of dimensions is smaller
than the ring dimension \cite{Roberts,RobertsRennes,GilletSoule}.
Their arguments use local Chern characters and Adams operations,
respectively. Gabber proved nonnegativity in the remaining
mixed-characteristic case using alterations and an intersection
calculation on a normal bundle
\cite[Section~6.1, Theorem~6.1.1]{Gabber}. Vanishing removes
lower-dimensional contributions in the reductions below. For
complementary-dimensional pairs, strict positivity must additionally
rule out zero.

Further progress in ramified mixed characteristic includes Skalit's
positivity theorem for regular local rings essentially smooth over a
two-dimensional regular local base \cite{Skalit}. KC and Soto Levins
prove positivity for modules satisfying a Serre-liftability condition
\cite[Theorem~A]{KCSotoLevins}. In their positivity statement,
$R=S/(f)$ for an unramified regular local ring $S$ and
$f\in\m_S\setminus\m_S^2$, and $M$ has a Serre lift $L$:
a finite $S$-module with $M=L\otimes_S R$ and
$\dim S-\dim L=\dim R-\dim M$; it need not satisfy the usual
Tor-vanishing condition for a lift. Theorem~\ref{thm:main} imposes
neither this liftability condition nor the smoothness hypothesis.

The proof combines normalized-length methods with algebraic $K$-theory.
Faltings introduced the normalized-length formalism in almost
mathematics \cite[Appendix~2]{Faltings}; Gabber--Ramero developed its
almost-ring-theoretic foundations further \cite{GabberRamero}.
The perfectoid length results of Cai--Lee--Ma--Schwede--Tucker
\cite{CLMST}, together with their characteristic-$p$ counterparts
\cite{Ma}, give algebras on which parameter Koszul complexes are exact
in positive degrees after passing to normalized length. The mixed-characteristic
construction uses perfectoidization and almost purity, as developed by
Bhatt--Scholze \cite{BhattScholze}. Land--Tamme's
descent theorem for truncating invariants \cite{LT} supplies the
categorical descent input. Bhatt--Hochster--Ma show that the existence
of lim Cohen--Macaulay sequences for both quotient domains implies
positivity \cite[Theorem~6.1]{BHM}. Ishiro--Shimomoto obtain such
sequences of algebras from perfectoid towers over complete
mixed-characteristic Noetherian local domains with $F$-finite residue field
\cite[Theorem~3.17]{IshiroShimomoto}. These are conditional routes
through finite modules. Here no such tower on either quotient domain
is assumed: the proof uses bounded complexes modulo free sequences
of negligible normalized rank.

\subsection{Proof strategy}

After standard reductions, it is enough to consider two complete local
domains $D=R/P$ and $E=R/Q$ of positive dimensions $h,l$, where
$h+l=d$ and $P+Q$ is $\m$-primary. Each domain has a finite extension,
denoted $D'$ or $E'$, with a finite Galois action over a regular
normalization base $A_D$ or $A_E$, that is, a regular local subring
over which the extension is finite. Associated algebras $C_D,C_E$ carry additive
nonnegative lengths $\lambda_D,\lambda_E$. Their parameter quotients
have positive length, whereas positive parameter Koszul homology has
length zero.

There are two distinct comparisons in the argument. First, rational
$K$-theory with support at the closed point shows that ordinary and
normalized Euler lengths agree after summing over all Galois
conjugates. Write $G_D,G_E$ for the two groups, and
$a_D=\rank_D D'$, $a_E=\rank_E E'$. For $\sigma\in G_D$, let
$c_\sigma$ be the normalized Tor Euler characteristic of $C_D$ and
$E'$ over $R$, where the $R$-action on $C_D$ is twisted by $\sigma$.
The first comparison identifies the sum of the $c_\sigma$ with
$|G_D|a_Da_E\chi^R(D,E)$. It therefore suffices to show that
each $c_\sigma$ is positive. Fix one such twist.

To analyze that value, choose a regular sequence $x$ in $Q$ that is a
parameter system on $D$. The module $U=C_D/xC_D$ has positive finite
normalized length over $B=R/(x)$. It can be approximated by finite-length
$B$-modules $M_n$ killed by one fixed power of the maximal ideal of $B$,
with ordinary lengths
divided by integers $s_n$ converging to the relevant normalized lengths.
Their minimal resolutions have ranks $o(s_n)$ above degree $l$.
Truncating at degree $l$ therefore produces supported complexes over
$E'$ in the category that kills free sequences whose ranks divided by
$s_n$ tend to zero.

The second comparison takes place in this category. Its supported
objects have nonnegative homology lengths, even when support appears
only after passing to the quotient. The two Euler evaluations again
agree on sums of conjugates. Evaluation using $C_E$ leaves only degree
zero, whose length is bounded below by a positive constant times the
normalized number of generators of $M_n$. The common annihilator and
the positive limiting length of $M_n$ make this lower bound strictly
positive. The second Galois average proves positivity for each value
in the first average, completing the argument.

Figure~\ref{fig:two-comparisons} records the two interfaces. Here
$b_{\sigma,\gamma}$ denotes the $C_E$-evaluation of the truncated
sequence after the additional twist $\gamma\in G_E$. Choose $a$ with
$\m_B^aM_n=0$ for all $n$, and write $\beta_0(M_n)$ for the minimal
number of generators of $M_n$. Then
$\length_B M_n\leq\beta_0(M_n)\length_B(B/\m_B^a)$.
Thus the positive normalized length of the $M_n$ forces a positive
normalized number of generators; it is this uniform estimate that
makes every $b_{\sigma,\gamma}$ strictly positive.

\begin{figure}[htbp]
\centering
\renewcommand{\arraystretch}{1.5}
\begin{tabular}{c}
\fbox{\strut Ordinary supported comparison over $A_D\subset D'$}\\
$\displaystyle |G_D|a_Da_E\,\chi^R(D,E)
       =\sum_{\sigma\in G_D}c_\sigma$\\
$\Big\uparrow$\quad prove positivity separately for each $\sigma$\\
\fbox{\strut Asymptotic supported comparison over $A_E\subset E'$}\\
$\displaystyle |G_E|c_\sigma
       =\sum_{\gamma\in G_E}b_{\sigma,\gamma}>0$
\end{tabular}
\caption{Two norm comparisons, not invariance of individual evaluations.
The first is Lemma~\ref{pos:first-average}; the second is
Equation~\eqref{pos:second-average}. The uniform annihilator and
generator bound give the strict inequality in
Lemma~\ref{pos:positive-evaluation}.}
\label{fig:two-comparisons}
\end{figure}

The reusable technical points are the finite-surjection descent
argument with arbitrary connective algebra coefficients and the
construction of supported Euler evaluations after quotienting by
negligible ranks. In particular, the comparison is proved on
idempotent summands in the quotient, not only on complexes that already
have uniform support before taking the quotient. No invariance of a
normalized-length function under either Galois group is assumed.

Section~\ref{sec:lengths} establishes the precise normalized-length
inputs. Section~\ref{sec:descent} proves supported rational descent and
the first Euler comparison. Section~\ref{sec:asymptotic} constructs the
rank quotient and its supported evaluations. Section~\ref{sec:positivity}
performs the approximation and the two comparisons to prove
Theorem~\ref{thm:main}.

\subsection{Consequences}

Theorem~\ref{thm:main} also has consequences for symbolic powers.
For a ramified mixed-characteristic regular local ring $(R,\m)$ and
primes $P,Q$ with $\sqrt{P+Q}=\m$ and
$\operatorname{ht}P+\operatorname{ht}Q=\dim R$, the implication of
Kurano--Roberts \cite[Theorem~3.2]{KuranoRoberts} now gives
\[
 P\cap Q^{(r)}\subseteq\m^{r+1}\qquad(r\geq1),
 \qquad Q^{(r)}:=Q^rR_Q\cap R.
\]
Here ramified means that the residue characteristic $p$ lies in
$\m^2$. This is a direct application of their positivity-to-symbolic-power
theorem, rather than an additional ingredient in our proof.

The next consequence applies Skalit's finite-support blowup argument
\cite[Corollary~D]{Skalit} with Theorem~\ref{thm:main} as its positivity
input. Dutta studied the lower bound in this finite-support setting
\cite{DuttaFiniteSupport}. We include the argument to exhibit how
strict positivity at the points of the blowup controls equality.

\begin{corollary}[Finite strict-transform intersection]\label{cor:tangent-cone}
Let $(A,\m)$ be a commutative Noetherian regular local ring with residue
field $k$, and let $P,Q$ be prime ideals such that $A/(P+Q)$ has finite
length and $\dim(A/P)+\dim(A/Q)=\dim A$. Write $e_\m$ for Hilbert--Samuel
multiplicity and put
\[
 T=\operatorname{gr}(A/P)\otimes_{\operatorname{gr}A}\operatorname{gr}(A/Q),
\]
where all associated graded rings use the maximal-ideal filtrations.
If $\dim T\leq1$, then
\[
 \chi^A(A/P,A/Q)\geq e_\m(A/P)e_\m(A/Q),
\]
with equality if and only if $\dim T=0$.
\end{corollary}

\begin{proof}
If $\dim A=0$, then $A=k$, $P=Q=0$, and both sides are $1$ with $T=k$.
Otherwise, let $\widetilde X$ be the blowup of $X=\Spec A$ at its closed
point, and let $\widetilde Y,\widetilde Z$ be the strict transforms of
$Y=\Spec(A/P)$ and $Z=\Spec(A/Q)$. By \cite[Lemma~3.5]{SkalitUnramified},
the underlying set $S$ of $\widetilde Y\cap\widetilde Z$ identifies with
that of $\operatorname{Proj}T$. Thus $S$ is a finite set of closed points
of the exceptional divisor, empty exactly when $\dim T=0$.
Fulton's blowup formula and Skalit's finite-support calculation
\cite[Theorem~3.4 and proof of Proposition~3.10]{SkalitUnramified} give
\[
 \chi^A(A/P,A/Q)-e_\m(A/P)e_\m(A/Q)
   =\sum_{x\in S}[\kappa(x):k]\,\chi^{B_x}(M_x,N_x),
\]
where $B_x=\mathcal O_{\widetilde X,x}$,
$M_x=\mathcal O_{\widetilde Y,x}$, and $N_x=\mathcal O_{\widetilde Z,x}$.
Each residue degree is finite and positive since the exceptional divisor
is projective over $k$. The blowup $\widetilde X$ is regular
\cite[Proposition~A.4(1)]{Jannsen}. The modules $M_x,N_x$ are nonzero and
finite over $B_x$, with finite-length tensor product because their
intersection is supported on $S$. The local domains $A/P$ and $A/Q$,
as well as $A$, are formally equidimensional and universally catenary.
Thus \cite[Lemma~3.8 and proof of Proposition~3.10]{SkalitUnramified}
gives $\dim M_x=\dim(A/P)$, $\dim N_x=\dim(A/Q)$, and
$\dim B_x=\dim A$. Theorem~\ref{thm:main} therefore makes every local
intersection multiplicity in the sum strictly positive. The asserted
inequality follows, and equality holds exactly when $S$ is empty.
\end{proof}

\subsection{Conventions}

We fix universes large enough for the module and sequence categories
under consideration.
Complexes are indexed homologically. An associative derived algebra is
\emph{connective} when its homotopy groups in negative degrees vanish.
Perfect modules over associative algebras are right modules.
All scalar localizations of coefficient algebras are along the central
image of the commutative base. We use nonconnective algebraic
$K$-theory of perfect complexes; a subscript $\Q$ denotes rationalization.
For a sequence $a=(a_1,\ldots,a_r)$, write $K(a;V)$ for its homological
Koszul complex with coefficients in $V$, and put
$H_i(a;V)=H_i(K(a;V))$. The notation $a^v$ denotes
$(a_1^v,\ldots,a_r^v)$, not the $v$th power of the ideal it generates.
Normalized lengths are allowed to vanish on nonzero modules; all uses
of their additivity involve finite values or nonnegative extended
values for which the indicated operations are defined.

\section{Residue fields and normalized length}\label{sec:lengths}

The comparison arguments require algebras on which positive parameter
Koszul homology has length zero and parameter quotients have prescribed
positive length.  The finite approximations in Section~\ref{sec:positivity}
also require that this length be defined on underlying modules over a
flat tower, without an action of the full algebra.  We construct these
objects after arranging a perfect residue field; neither unramifiedness
nor Cohen--Macaulayness of the domain is required.
Both characteristics are needed: a prime quotient of the
mixed-characteristic ambient ring may itself have characteristic $p$,
even though ambient equicharacteristic positivity is already known.

\subsection{Extending the residue field}

\begin{lemma}[Residue-field extension]\label{len:residue-extension}
Let $(B,\mathfrak b,k)$ be a complete Noetherian local ring, and let
$\overline{k}$ be an algebraic closure of $k$.  There is a complete
Noetherian local ring $(B',\mathfrak b',\overline{k})$ and a flat local map
$B\to B'$ such that $\mathfrak bB'=\mathfrak b'$.
For every finite $B$-module $V$ and every $a\geq1$,
\[
 \length_{B'}\bigl((V\otimes_B B')/(\mathfrak b')^a
                          (V\otimes_B B')\bigr)
       =\length_B(V/\mathfrak b^aV).
\]
Consequently module dimensions are preserved, as are finite lengths.
If $B$ is regular, then $B'$ is regular.  For finite $B$-modules $M,N$
whose tensor product has finite length, every Tor length, and hence
their intersection multiplicity when defined, is preserved.
\end{lemma}

\begin{proof}
In equal characteristic, a coefficient field gives a presentation
$B=k[[X_1,\ldots,X_t]]/I$.  Put
$B'=\overline{k}[[X_1,\ldots,X_t]]/I\overline{k}[[X_1,\ldots,X_t]]$.
The map of power-series rings is flat, by the local flatness criterion
for the variables and the flat coefficient-field extension.  Its base
change to the quotient has the asserted residue field and maximal ideal.

In the remaining case the residue characteristic is a prime $p$.
Cohen structure supplies a presentation $B=V[[X_1,\ldots,X_t]]/I$,
where $V$ is a complete discrete valuation ring with uniformizer $p$
and residue field $k$.  We construct a discrete valuation extension
with unchanged uniformizer and residue field $\overline{k}$, without
assuming that $k$ is perfect.  Given a simple algebraic residue-field
extension, lift its monic irreducible polynomial to $f\in V[X]$.
The algebra $V[X]/(f)$ is finite free over $V$, its reduction modulo $p$
is the desired field, and all its maximal ideals lie over $(p)$.
It is therefore local with maximal ideal generated by $p$.
Since $p$ is a nonzerodivisor, it is a one-dimensional regular local
ring, hence a discrete valuation ring.  This argument also applies
to inseparable residue extensions.

Perform this construction along a tower of simple algebraic extensions
exhausting $\overline{k}$, taking unions at limit stages.  In every such
union, a nonzero element is a unit times a nonnegative integral power
of $p$.  Every nonzero ideal is generated by an element of least
valuation, so the union is still a discrete valuation ring.  Its
completion $V'$ has uniformizer $p$ and residue field $\overline{k}$.
The map $V\to V'$ is flat because $V'$ is torsion-free over $V$.
Now take
\[
 B'=V'[[X_1,\ldots,X_t]]/IV'[[X_1,\ldots,X_t]].
\]
The same local flatness criterion proves the required flatness.
This is the usual residue-extension construction; compare
\cite[Lemma~2.2]{Ma}.

For either construction, tensor a composition series of a finite-length
$B$-module with $B'$.  Flatness preserves exactness, and every simple
factor becomes $B'/\mathfrak b'$.  This proves the length assertions
and, applied to $V/\mathfrak b^aV$, the displayed identity.  Equality of
Hilbert--Samuel functions gives equality of module dimensions.  If $B$
is regular of dimension $d$, its $d$ maximal-ideal generators generate
$\mathfrak b'$, and $\dim B'=d$, so $B'$ is regular.  Finally flat base
change identifies the Tor modules after extension; their finite
lengths are preserved by the preceding argument.
\end{proof}

\subsection{Regular bases and the tower length}

For a perfect field $k$ of characteristic $p$, let $W(k)$ denote its
ring of Witt vectors.

\begin{lemma}[Regular normalization bases]\label{len:normalization}
Let $(D,\mathfrak d,k)$ be a complete Noetherian local domain of
dimension $h>0$, where $k$ is perfect of characteristic $p>0$.
There is a finite, generically separable inclusion $A\subseteq D$,
inducing the identity on residue fields, with
\[
 A=k[[t_1,\ldots,t_h]]\quad\text{if }\operatorname{char}D=p,
 \qquad
 A=W(k)[[t_2,\ldots,t_h]]\quad\text{if }\operatorname{char}D=0.
\]
In the second case set $t_1=p$.  In both cases $t_1,\ldots,t_h$
is a parameter system of $D$ and a regular parameter system of $A$.
\end{lemma}

\begin{proof}
The positive-characteristic assertion is the Cohen--Gabber
normalization theorem \cite{CohenGabber}.  In mixed characteristic,
choose a coefficient ring $W(k)\subseteq D$ and extend $p$ to a
parameter system $p,t_2,\ldots,t_h$.  Completeness and the parameter
condition make the resulting power-series map module finite.
Its kernel is zero: the source and target have dimension $h$, whereas
a nonzero ideal of the regular local domain
$W(k)[[t_2,\ldots,t_h]]$ lowers dimension.  The fraction-field
extension is separable in characteristic zero.
\end{proof}

Fix one of the bases in Lemma~\ref{len:normalization} and compatible
$p$-power roots of its parameters.  Write
\begin{equation}\label{len:tower-definition}
 A_n=A[t_1^{1/p^n},\ldots,t_h^{1/p^n}],\qquad
 \Lambda=\varinjlim_n A_n.
\end{equation}
The tower is uncompleted.  Only in mixed characteristic do we also
use its $p$-adic completion, denoted $A_\infty$.

\begin{lemma}[The tower and its length]\label{len:tower}
Each $A_n$ is a complete regular local ring with residue field $k$.
For $b\geq a$, the map $A_a\to A_b$ is finite free of rank
$p^{(b-a)h}$; in particular $A\to A_n$ has rank $p^{nh}$.
On the category of $\Lambda$-modules killed by some power of
$\mathfrak m_A$ there is a nonnegative additive function
$\lambda_\infty$ with the following properties.
\begin{enumerate}
\item A finitely presented module descends to a finite-length
      $A_n$-module $V_n$, and
      \begin{equation}\label{len:finite-stage}
       \lambda_\infty(V_n\otimes_{A_n}\Lambda)
                      =p^{-nh}\length_{A_n}(V_n).
      \end{equation}
\item For a finitely generated module $V$, its length is the infimum
      of the lengths of finitely presented modules in this category
      admitting a surjection onto $V$.
\item For an arbitrary module $V$ in this category,
      \begin{equation}\label{len:supremum}
       \lambda_\infty(V)=
       \sup\{\lambda_\infty(V_0):V_0\subseteq V
                                  \text{ is finitely generated}\}.
      \end{equation}
\end{enumerate}
The value on an arbitrary module may be infinite.  In mixed
characteristic, restriction from $A_\infty$ to $\Lambda$ gives the same
length on modules killed by a fixed power of $\mathfrak m_A$.
\end{lemma}

\begin{proof}
In characteristic $p$, the stages are power-series rings in the chosen
roots.  In mixed characteristic they are power-series rings over the
complete discrete valuation ring $W(k)[p^{1/p^n}]$.
The monomials in the successive roots, with each exponent less than
the relevant root degree, form bases for all transition extensions.
This proves the assertions about regularity, ranks, and residue fields.

Finite presentations over the union descend to a stage.  The additional
relations expressing annihilation by a fixed power of
$\mathfrak m_A$ also descend after increasing that stage, because
there are finitely many generators and relations to consider.  The
resulting finite-stage module has finite length.  Formula
\eqref{len:finite-stage} is unchanged on enlarging the stage, by
flatness and the computed free ranks.  The infimum and supremum
extensions, including their additivity and compatibility with this
formula, are Faltings' normalized-length formalism
\cite[Appendix~2]{Faltings}, in the formulation of
\cite[Definition~2.2.1 and Propositions~2.2.3, 2.2.6(a)]{CLMST}.
The same residue field implies that $A_n$-length and $A$-length
agree on finite-length $A_n$-modules.

In mixed characteristic $\Lambda$ is $p$-torsion-free, and the natural
map induces $\Lambda/p^a\Lambda\simeq A_\infty/p^aA_\infty$ for every
$a\geq1$.  A module killed by $\mathfrak m_A^a$ is killed by $p^a$.
Its submodules and finite generating sets are therefore identical
over the two bases.  In computing the infimum over finitely presented
covers of such a module, one may replace a cover $V$ by
$V/\mathfrak m_A^aV$; this is still finitely presented.  These covers
also correspond over the identical quotient rings, proving the last
assertion.
\end{proof}

\begin{lemma}[Zero-length limits]\label{len:zero-length-limits}
Suppose all modules under consideration are killed by one fixed power
of $\mathfrak m_A$.  Sums, direct sums, and directed colimits of
zero-length modules have length zero.  For an increasing union,
\[
 \lambda_\infty\Bigl(\bigcup_a V_a\Bigr)
                         =\sup_a\lambda_\infty(V_a).
\]
If $(V_a)_{a\geq1}$ is a direct system with
$\lambda_\infty(V_a)\to0$, then
$\lambda_\infty(\varinjlim_a V_a)=0$, without any injectivity
assumption on the transition maps.
\end{lemma}

\begin{proof}
A finite set of elements of a direct sum is contained in a finite
subsum.  Additivity and \eqref{len:supremum} give the assertion for
direct sums; sums and directed colimits are quotients of direct sums.
For an increasing union, every finite generating set lies in one
stage.  Finally a finitely generated submodule of $\varinjlim_a V_a$
is contained in the image of every sufficiently late $V_a$, by lifting
its generators to a common stage.  Monotonicity bounds its length by
$\lambda_\infty(V_a)$ at all such stages.  Apply
\eqref{len:supremum}.  The common annihilator ensures throughout
that these modules belong to the length category.
\end{proof}

\subsection{Length algebras and parameter homology}

The tower length will now measure modules over an algebra of the given
domain.  Its parameter identities have two uses: the base parameters
control the Euler comparisons, while arbitrary parameter systems give
the positive quotients used in Section~\ref{sec:positivity}.
For an ideal $I$ primary to the maximal ideal of a local ring $D$, write
$e(I,D)$ for its Hilbert--Samuel multiplicity.

Related normalized-length constructions, in which parameter-quotient
lengths recover multiplicities and positive parameter Koszul homology has
normalized length zero, appear in \cite[Theorem~3.1]{OpenAILech2026}.
Here we retain the chosen base $A\subseteq D$ and the finite-stage,
infimum, and supremum properties on underlying tower modules.

\begin{theorem}[Normalized-length input]\label{len:input}
Let $(D,\mathfrak d,k)$ and a chosen finite generically separable base
$A\subseteq D$ be as in Lemma~\ref{len:normalization}, and let
$\Lambda$ be the tower \eqref{len:tower-definition}.  There is a
$D$-algebra $C_D$, carrying a compatible $\Lambda$-action, such that
restriction of scalars defines an additive nonnegative length
$\lambda_D=\lambda_\infty$ on its modules annihilated by a power of
$\mathfrak d$.  For every parameter system $z_1,\ldots,z_h$ of $D$
and every integer $v\geq1$,
\begin{align}
 \lambda_D H_i(z_1^v,\ldots,z_h^v;C_D)&=0\qquad(i>0),
                                      \label{len:koszul-vanishing}\\
 \lambda_D\bigl(C_D/(z_1^v,\ldots,z_h^v)C_D\bigr)
                    &=v^h e((z_1,\ldots,z_h),D).
                                      \label{len:parameter-length}
\end{align}
The right side is positive and finite.  The algebra and length are
independent of the parameter system $z$.
For the regular parameters of $A$, in particular,
\begin{equation}\label{len:base-parameter-length}
 \lambda_D\bigl(C_D/(t_1^v,\ldots,t_h^v)C_D\bigr)
                         =v^h\rank_A D.
\end{equation}
The finite-stage, infimum, and supremum properties of
Lemma~\ref{len:tower} remain available on underlying tower modules;
they do not require that the relevant submodules or covers be
$C_D$-modules.
\end{theorem}

\begin{proof}
We give the positive-characteristic construction and then specify the
mixed-characteristic literature inputs.

Assume first that $D$ has characteristic $p$, and put
$C_D=D_{\mathrm{perf}}=\bigcup_n D^{1/p^n}$ inside a fixed algebraic
closure of its fraction field. Here $\Lambda=A_{\mathrm{perf}}$.
Write $K=\operatorname{Frac}(A)$, $L=\operatorname{Frac}(D)$, and
$K'=K_{\mathrm{perf}}$.  For $q=p^n$ put
\[
 B_q=\Lambda\otimes_{A^{1/q}}D^{1/q}.
\]
Flatness of $\Lambda$ over $A^{1/q}$ and torsion-freeness of
$D^{1/q}$ embed this tensor product into
$K'\otimes_{K^{1/q}}L^{1/q}$.  Generic separability makes the latter
a field, by linear disjointness from a purely inseparable extension.
Thus $B_q=\Lambda D^{1/q}\subseteq C_D$, and these subrings increase
to $C_D$.

Let $H$ be any $\mathfrak d$-primary ideal and set
$V_q=B_q/HB_q$ and $W_q=\operatorname{im}(V_q\to C_D/HC_D)$.
Write $H^{[q]}=(a^q:a\in H)$ for its Frobenius power.
Finite-stage normalization and the $q$-th-power isomorphism give
\[
 \lambda_\infty(V_q)=q^{-h}\length_D(D/H^{[q]}).
\]
Although $B_q\subseteq C_D$, the map on these quotients need not be
injective.  To recover the length of $C_D/HC_D$ from the displayed
formula, we control the kernels of $V_q\to W_q$ by a single trace
denominator.  There is a nonzero $g\in A$ with
\begin{equation}\label{len:conductor}
 gC_D\subseteq B_1,\qquad g^{1/q}C_D\subseteq B_q.
\end{equation}
Indeed, choose a $K$-basis $b_1,\ldots,b_r$ of $L$ in $D$ and its
trace-dual basis $\beta_1,\ldots,\beta_r$.  A common denominator
$g\in A\setminus\{0\}$ makes every $g\beta_j$ lie in $D$.
The field $LK'$ is perfect and equals $L_{\mathrm{perf}}$.
The ring $\Lambda$ is normal, since an integral equation and the
fractions occurring in it descend to a normal stage $A^{1/q}$.
For $c\in C_D$, each
$\operatorname{Tr}_{LK'/K'}(g\beta_jc)$ is integral over $\Lambda$
and belongs to $K'$, hence belongs to $\Lambda$.  The trace-dual
expansion of $gc$ proves the first inclusion in
\eqref{len:conductor}.  Taking $q$-th roots proves the second.

The kernel of $V_q\to W_q$ is killed by $g^{1/q}$, by
\eqref{len:conductor}.  Kernel and cokernel of an endomorphism of a
finite-length object have equal length.  Applied to multiplication
by $g^{1/q}$, this yields
\[
 \lambda_\infty\ker(V_q\to W_q)
 \leq q^{-h}\length_D D/(H^{[q]},g)=O(q^{-1}).
\]
For the estimate, if $H$ has $s$ generators then
$H^{s(q-1)+1}\subseteq H^{[q]}$.  Thus $H^{[q]}$ contains a power
of $\mathfrak d$ with exponent $O(q)$, whereas
$\dim D/(g)\leq h-1$ if $g$ is a nonunit.  Hilbert--Samuel growth
gives the asserted bound; for a unit $g$ the quotient is zero.
The $W_q$ increase to $C_D/HC_D$ and are uniformly primary-torsion.
The Hilbert--Kunz limit and Lemma~\ref{len:zero-length-limits} show
\begin{equation}\label{len:perfection-quotient}
 \lambda_D(C_D/HC_D)=e_{\mathrm{HK}}(H,D).
\end{equation}
Here we use Monsky's existence theorem for Hilbert--Kunz multiplicity
\cite{Monsky}; see also \cite[Section~2.2]{Ma}.

We have identified primary-quotient lengths with Hilbert--Kunz
multiplicities.  For a parameter ideal, it remains to obtain the
Hilbert--Samuel value and zero positive Koszul homology.  If $z$ is
a parameter system of $D$, finite-stage flatness gives
\[
 \lambda_\infty H_i(z;B_q)
          =q^{-h}\length_D H_i(z_1^q,\ldots,z_h^q;D).
\]
The Koszul complex $K(z;D)$ has finite free terms in degrees
$0,\ldots,h$ and finite-length homology.  The Frobenius homology
theorem of Roberts, in the form \cite[Theorem~2.4]{Ma}, says that
the displayed quantity tends to zero for $i>0$; the original result
is \cite{RobertsIntersectionTheorems}. Its hypotheses
require precisely a complete characteristic-$p$ ring and a
finite-free complex of length its dimension with finite-length
homology.  Filtered colimits commute with homology, and these
Koszul homology modules have the common annihilator $(z)$.
Lemma~\ref{len:zero-length-limits} proves
\eqref{len:koszul-vanishing}, also with $z$ replaced by $z^v$.
Finally the Koszul Euler characteristic formula gives
\[
 \sum_{i=0}^h(-1)^i\length_D H_i(z_1^q,\ldots,z_h^q;D)
                           =q^h e((z),D).
\]
After division by $q^h$, the positive-degree terms tend to zero.
Together with \eqref{len:perfection-quotient} and the multiplicity
power rule this proves \eqref{len:parameter-length}.

Now suppose $D$ has mixed characteristic.  Let $A_\infty$ be the
$p$-adic completion of $\Lambda$ and define
\[
 C_D=(D\otimes_A A_\infty)_{\mathrm{perfd}},
\]
the perfectoidization of \cite[Notation~4.0.1]{CLMST}, built on
Bhatt--Scholze's construction \cite[Proposition~8.5]{BhattScholze}.
The stated hypotheses there are exactly that $D$ is a complete
mixed-characteristic local domain with perfect residue field and
the chosen finite Witt power-series base $A$.
Choose $g\in A\setminus\{0\}$ so that $A[1/g]\to D[1/g]$ is
finite etale.  Corollary~4.0.4 of \cite{CLMST} annihilates positive
Koszul homology for every parameter system of $A_n\otimes_A D$
by the almost ideal $(g)_{\mathrm{perfd}}$.
At $n=0$ this includes every parameter system of $D$.
Proposition~4.0.10 of that paper makes these uniformly
$\mathfrak m_A$-power-torsion homology modules have normalized
length zero.  Proposition~4.0.14, again at $n=0$, identifies the
length of the parameter quotient with its Hilbert--Samuel
multiplicity.  These two results, applied also to $z^v$, prove
\eqref{len:koszul-vanishing} and \eqref{len:parameter-length}
for this single algebra $C_D$.  Lemma~\ref{len:tower} permits
restriction to the uncompleted tower throughout.

In either characteristic, a module killed by a power of
$\mathfrak d$ is killed by a power of $\mathfrak m_A$, so all
claimed lengths are defined.  Additivity and the tower properties
are those already established.  Lastly the multiplicity rank formula,
and equality of the residue fields, give
$e((t),D)=\rank_A D\,e((t),A)=\rank_A D$.
This proves \eqref{len:base-parameter-length}.
\end{proof}

The base-parameter case extends to every finite-length module over
$A$.  The same parameter-quotient positivity also prevents the residue
quotient of $C_D$ from having length zero; this will supply the positive
residue-length factor in the final degree-zero estimate.

\begin{corollary}[Finite-length base change]\label{len:finite-length-base-change}
In Theorem~\ref{len:input}, put $r=\rank_A D$.
For every finite-length $A$-module $V$,
\[
 \lambda_D\Tor_i^A(V,C_D)=0\quad(i>0),\qquad
 \lambda_D(V\otimes_A C_D)=r\length_A(V).
\]
Moreover $\lambda_D(C_D/\mathfrak d C_D)>0$.
\end{corollary}

\begin{proof}
For $V=k$, the first two assertions are
\eqref{len:koszul-vanishing} and \eqref{len:base-parameter-length},
using the Koszul resolution over regular $A$.
Induct on a composition series of $V$.  The Tor long exact
sequences and length additivity discard all positive-degree terms
and show that degree-zero length increases by $r$ at each step.
All modules here have bounded primary torsion.

For the last assertion choose a parameter ideal $H$ of $D$.
Tensor a composition series of $D/H$ with $C_D$.
Right exactness gives a filtration of $C_D/HC_D$ whose factors
are quotients of $C_D/\mathfrak dC_D$.  Therefore
\[
 0<e(H,D)=\lambda_D(C_D/HC_D)
       \leq\length_D(D/H)\lambda_D(C_D/\mathfrak dC_D).
\]
This proves the required strict positivity.
\end{proof}

\section{Rational descent with closed supports}
\label{sec:descent}

We compare ordinary and normalized Euler lengths through rational
$K$-theory with support. In the regular-base applications, the strategy
is to compare the two evaluations on classes pulled back from the base,
then show that sums of conjugate classes come from rational classes over
that base. The latter step is the descent theorem below; the final
subsection combines it with the normalized-length input of
Section~\ref{sec:lengths}. We retain connective associative coefficients
because Section~\ref{sec:asymptotic} applies the same descent theorem to
the quotient by negligible-rank sequences.

All derived algebras in this section are indexed homologically. Thus
connective means that their homotopy groups vanish in negative degrees.

Let $A$ be a Noetherian commutative $\Z_{(p)}$-algebra of finite Krull
dimension, let $J\subset A$ be an ideal containing $p$, and let
$H$ be a connective associative $A$-algebra. The image of $A$ acts
centrally. For a commutative $A$-algebra $B$, write
\[
 \Perf_J(H\otimes_A^{\mathbf L}B)
\]
for the category of perfect modules whose central localizations vanish
off $V(JB)$. Its nonconnective $K$-theory spectrum and its zeroth
homotopy group are denoted by
$\K^J(H\otimes_A^{\mathbf L}B)$ and
$K_0^J(H\otimes_A^{\mathbf L}B)$, respectively. Rationalization of a
spectrum is indicated by a subscript $\Q$.
For a finite group action, $(-)^{hG}$ denotes homotopy fixed points
(derived invariants), whereas $(-)^G$ denotes ordinary invariants.

\begin{theorem}[Supported rational descent]
\label{ds:descent}
Let $A$ be a commutative Noetherian $\Z_{(p)}$-algebra of finite Krull
dimension, let $J\subset A$ be an ideal containing $p$, and let $H$ be a connective
associative $A$-algebra with central $A$-action.
Let $T$ be a finite commutative $A$-algebra with an action of a finite
group $G$ over $A$. Suppose that $\Spec T\to\Spec A$ is surjective and
that $G$ acts transitively on the points of every geometric fiber.
Then pullback induces an equivalence
\begin{equation}
 \K^J(H)_\Q
 \simeq
 \bigl(\K^J(H\otimes_A^{\mathbf L}T)_\Q\bigr)^{hG}.
 \label{ds:spectral-descent}
\end{equation}
In particular, the map
\begin{equation}
 K_0^J(H)\otimes\Q
 \longrightarrow
 \bigl(K_0^J(H\otimes_A^{\mathbf L}T)\otimes\Q\bigr)^G
 \label{ds:kzero-descent}
\end{equation}
is surjective.
\end{theorem}

Neither flatness nor tame ramification of $A\to T$ is required.
The coefficient algebra $H$ need not be commutative, finite over $A$,
or bounded in positive homological degrees. Its connectivity is
essential to the argument.

\subsection{The truncating invariant and its supports}

We use the following precise form of the results of Land--Tamme.
A localizing invariant takes exact sequences of stable categories to
fiber sequences. It is called \emph{truncating} if its value on a
connective associative algebra is unchanged by passage to $\pi_0$.
Write $\operatorname{fib}$ for homotopy fiber and
$\operatorname{HN}(-/\Q)$ for negative cyclic homology over $\Q$.
Rational infinitesimal $K$-theory
\[
 I(U)=\operatorname{fib}\bigl(
   \K(U)_\Q\longrightarrow
   \operatorname{HN}(U\otimes\Q/\Q)\bigr)
\]
is truncating by Goodwillie's relative comparison theorem
\cite[Main Theorem]{Goodwillie}, with the passage from simplicial rings
to connective associative ring spectra explained in
\cite[proof of Corollary~3.9 and Remark~3.10]{LT}. Truncating localizing
invariants satisfy cdh descent, also after tensoring with a fixed
connective associative algebra over the base
\cite[Theorem~A.2 and Appendix~B, especially Theorem~B.3]{LT}.
The cdh topology is generated by Nisnevich covers and abstract blowup
squares: a proper map and a closed center, with the proper map an
isomorphism off that center.

For a derived $\Z_{(p)}$-algebra $U$, consider
\begin{equation}
 E(U)=\operatorname{fib}\bigl(
      \K(U)_\Q\longrightarrow\K(U[1/p])_\Q\bigr).
 \label{ds:relative-invariant}
\end{equation}
The negative cyclic homology terms in the two character maps are
naturally identical, since
$U\otimes\Q\simeq U[1/p]\otimes\Q$. Taking fibers of the character
square therefore gives
\begin{equation}
 E(U)\simeq\operatorname{fib}\bigl(I(U)\longrightarrow I(U[1/p])\bigr).
 \label{ds:character-cancellation}
\end{equation}
It follows that $E$ is a truncating localizing invariant of
$\Z_{(p)}$-linear categories. Notice that this argument does not
identify $\K(U)_\Q$ with $\K(U\otimes\Q)$.

Tensoring with $H$ preserves the localizing property. It also
preserves truncation: if $B$ is connective, then
$H\otimes_A^{\mathbf L}B$ is connective and
\[
 \pi_0(H\otimes_A^{\mathbf L}B)
   =\pi_0(H)\otimes_A\pi_0(B).
\]
Thus replacing $B$ by $\pi_0(B)$ does not change the value of $E$.
For a finite-type $A$-scheme $Y$, let $H_Y$ be the derived pullback
of $H$ and put $Z_Y=Y\times_A V(J)$. Define
\begin{equation}
 \cL(Y)=\operatorname{fib}\bigl(
    E(Y;H)\longrightarrow E(Y\setminus Z_Y;H)\bigr).
 \label{ds:supported-functor}
\end{equation}
Here the coefficient notation means evaluation on the category of
perfect $H_Y$-modules. The coefficient theorem just cited gives cdh
descent for $E(-;H)$. Restriction to the open complement of $Z_Y$
preserves the squares defining this topology, and limits commute
with limits. Hence $\cL$ also satisfies cdh descent.

Because $p\in J$, the support $Z_Y$ disappears after inverting $p$.
Expanding the two fibers in \eqref{ds:supported-functor} and using
perfect-complex localization consequently identifies
\begin{equation}
 \cL(Y)\simeq\K\bigl(\Perf_{Z_Y}(H_Y)\bigr)_\Q.
 \label{ds:support-identification}
\end{equation}
In particular, this is the desired rational supported $K$-theory,
not a cdh replacement of it. The functor is nilinvariant as well.

We recall why the support condition is compatible with the
coefficient categories used here. If $P$ is perfect and $z\in A$,
compactness identifies localization of its endomorphism spectrum
with the endomorphism spectrum of $P[1/z]$. Therefore
$P[1/z]=0$ if and only if some power $z^v$ kills
$\operatorname{id}_P$ in the homotopy category. Applying this to
finitely many generators of $J$ gives the central support
criterion. The corresponding Koszul tensors generate the compact
kernel of restriction to the complement. Perfect localization,
including idempotent completion, thus computes the fiber in
\eqref{ds:support-identification}.

The geometric localization input is Thomason--Trobaugh's perfect
localization theorem \cite[Propositions~5.2.2--5.2.4 and
Theorem~7.4]{ThomasonTrobaugh}. Tensoring its exact sequence with the
perfect coefficient category, using \cite[Lemmas~3.12 and B.2]{LT}
and the preceding supported-kernel identification, gives the
coefficient version. Nonconnective $K$-theory sends this exact sequence
to a fiber sequence by \cite[Definition~8.1, Theorems~8.7 and 9.8,
and Section~9.3]{BGT}.

\subsection{Descent for finite surjections}

The functor $\cL$ is now a nilinvariant cdh sheaf computing actual
rational supported $K$-theory. To apply it to the finite map in
Theorem~\ref{ds:descent}, we need descent for finite surjections, which
need not be flat. Rational transfers give the finite flat case;
flattening by blowup and induction on closed subsets will give the
general case.

\begin{lemma}
\label{ds:finite-flat}
The functor $\cL$ satisfies ordinary \v{C}ech descent for finite
faithfully flat covers, universally under base change.
\end{lemma}

\begin{proof}
Let $f:X\to Y$ be finite faithfully flat. Since the schemes are
Noetherian, $f_*\mathcal O_X$ is a vector bundle of positive rank.
Restriction of scalars preserves perfect coefficient modules and
their supports: locally the scalar-extended coefficient algebra is
a finite projective module over the original one. It therefore
defines a transfer. Finite-flat base change and the projection
formula give natural transformations of presheaves on $Y$-schemes
\[
 \cL\xrightarrow{f^*}\cL_X
       \xrightarrow{\operatorname{tr}_f}\cL,
 \qquad
 \cL_X(W)=\cL(W\times_Y X),
\]
whose composite is tensoring with the pullback of
$f_*\mathcal O_X$.

Work first on a Zariski open of $Y$ on which that vector bundle is
free of rank $a>0$, and fix a trivialization there. Additivity of
$K$-theory identifies this composite with multiplication by $a$.
The same trivialization pulls back on every $Y$-scheme, so this
identity is natural on the whole presheaf diagram. Rationally,
$a^{-1}\operatorname{tr}_f$ makes $\cL$ a retract of $\cL_X$.

After base change to $X$, the augmented \v{C}ech nerve of $f$ has
a section and hence an extra degeneracy. Its image under any
presheaf has split descent. Thus $\cL_X$ satisfies descent for
$f$. The descent comparison for $\cL$ is a retract of this
equivalence, and is itself an equivalence. Finally, both its source
and its totalization are Zariski sheaves. The conclusion therefore
glues over the trivializing cover of $Y$. All constructions commute
with base change.
\end{proof}

\begin{lemma}
\label{ds:finite-surjection}
The functor $\cL$ satisfies ordinary \v{C}ech descent for every
finite surjection of finite-type $A$-schemes.
\end{lemma}

\begin{proof}
Let $\tau$ be the topology generated by cdh covers and finite
faithfully flat covers. These cover families are stable under base
change. By Lemma~\ref{ds:finite-flat} and cdh descent, $\cL$ is a
$\tau$-sheaf. We show that every finite surjection $X\to Y$ is a
$\tau$-cover.

Reductions and finite closed decompositions are cdh-local
operations. For clarity, if a reduced scheme is a union
$Y=Y_1\cup Y_2$ of reduced closed subschemes, the proper map
$Y_1\to Y$ is an isomorphism off $Y_2$ and gives the abstract
blowup square with intersection $Y_1\cap Y_2$. This reduces closed
unions to their pieces. We may consequently reduce to integral
$Y$, with induction on dimension for proper closed subsets.

Choose a reduced irreducible component $X_0$ of $X$ dominating
$Y$. The map $X_0\to Y$ is finite and surjective. On a dense open
$U\subset Y$ it is finite flat of positive rank. Flattening by
blowup, applied to $\mathcal O_{X_0}$, gives a blowup
$Y'\to Y$, an isomorphism over $U$, for which the strict transform
$X'_0\to Y'$ is flat and finitely presented
\cite[Tag~0815]{Stacks}. The strict transform is a closed subscheme
of $X_0\times_Y Y'$, so it is finite over $Y'$. The blowup $Y'$ is
integral, and the image of $X'_0$ contains its generic point.
Finiteness makes that image closed. Thus $X'_0\to Y'$ is
surjective, and is a finite faithfully flat cover factoring through
$X\times_Y Y'$.

Let $Z\subsetneq Y$ be the blowup center. The family
$\{Y'\to Y,Z\to Y\}$ is a cdh cover. By induction on dimension,
the finite surjection $X\times_Y Z\to Z$ is a $\tau$-cover.
Together with $X'_0\to Y'$, these maps give a $\tau$-covering
family of $Y$ that factors through $X$. Hence $X\to Y$ itself
generates a covering sieve. The induction starts in dimension zero,
where a reduced integral base is a field and any nonzero finite
algebra is finite faithfully flat. This proves the assertion, and
therefore the required \v{C}ech descent.
\end{proof}

The argument uses actual retracts and sheaf limits; it requires
neither hyperdescent nor convergence of an unbounded descent
spectral sequence.

\subsection{The geometric orbit argument}

We have established descent for finite surjections. It remains to use
geometric transitivity to identify the homotopy $G$-invariants with the
value on the base. After arranging a section and reducing the schemes,
we compare its translates away from their coincidence loci and use
induction along those loci.

\begin{proof}[Proof of Theorem~\ref{ds:descent}]
Write $X=\Spec T$ and $S=\Spec A$. On finite-type $S$-schemes
consider the natural comparison
\begin{equation}
 \cL(Y)\longrightarrow
 \cL(X\times_S Y)^{hG}.
 \label{ds:comparison}
\end{equation}
Both sides are sheaves for the topology used in
Lemma~\ref{ds:finite-surjection}: base change preserves the cover
families, and homotopy fixed points commute with limits. Applying
finite-surjection descent to $X\to S$, it suffices to prove
\eqref{ds:comparison} on the nonaugmented terms of its \v{C}ech
nerve. Over each such term, the pulled-back $X$-cover admits a
section, given by a projection to one of its $X$-coordinates. The
geometric transitivity hypothesis is stable under base change.
We are therefore reduced to covers admitting a section.

We prove this case by induction on the dimension of $Y$.
Nilinvariance and the closed-union squares described in the proof
of Lemma~\ref{ds:finite-surjection} reduce the base to an integral
reduced scheme. Indeed both sides of \eqref{ds:comparison} have
these descent properties; the component reduction uses finitely
many closed pieces and lower-dimensional intersections.

Let $s:Y\to X\times_S Y$ be a section. Its finitely many
$G$-translates cover the source on points. To see this, extend the
residue field at any point of $Y$ to an algebraic closure and use
transitivity on that geometric fiber. Each translate is a closed
immersion, because the source is separated over $Y$, and each
factors through the reduction
\[
 V=(X\times_S Y)_{\mathrm{red}}.
\]
Retain the distinct translated sections and index them by a
transitive finite $G$-set $I$. Two distinct sections have a proper
closed coincidence locus in $Y$: their equalizer is closed, and
agreement at the generic point of integral reduced $Y$ would imply
agreement everywhere. Let $Z\subsetneq Y$ be the reduced union
of these coincidence loci, with $Z=\varnothing$ if there is only
one section.

The map
\[
 q:W=\coprod_{i\in I}Y\longrightarrow V
\]
is finite and surjective. Over $Y\setminus Z$ the section images
are pairwise disjoint and cover the reduced scheme $V$. Their
disjoint union is therefore isomorphic to $V$ there. Consequently
$q$ is a proper modification and the square
\[
 \begin{array}{ccc}
 W\times_Y Z&\longrightarrow&W\\
 \big\downarrow&&\big\downarrow\\
 V\times_Y Z&\longrightarrow&V
 \end{array}
\]
is an abstract blowup square. It is $G$-equivariant.

The $G$-action on $W$ merely permutes its copies of $Y$.
The canonical restriction maps give a $G$-equivariant equivalence
$\cL(W)\simeq\prod_{i\in I}\cL(Y)$.
A stabilizer of one section acts identically on that copy, including
on the coefficient algebra pulled back from $Y$. This is a
canonical identity of coefficient-bearing $Y$-schemes, so the
action is the ordinary permutation action, with coherently trivial
stabilizer action. The same statement holds after restriction to $Z$.
Since $\cL$ is rational, averaging over a finite group is exact.
It follows that
\begin{equation}
 \cL(W)^{hG}\simeq\cL(Y),
 \qquad
 \cL(W\times_Y Z)^{hG}\simeq\cL(Z).
 \label{ds:permutation-invariants}
\end{equation}
For example, if $I=G/G_s$, the first fixed-point spectrum is
$\cL(Y)^{hG_s}$ for the trivial $G_s$-action, and is $\cL(Y)$.
This identity holds without boundedness: the trivial
$\Q[G_s]$-module $\Q$ is projective, so derived invariants are
ordinary invariants even on unbounded rational complexes.

Apply $\cL$ to the abstract blowup square and take homotopy
$G$-invariants. The inductive hypothesis over $Z$ identifies
\[
 \cL(V\times_Y Z)^{hG}\simeq\cL(Z).
\]
Here nilinvariance allows one to pass between $V\times_Y Z$ and
the corresponding reduced or unreduced pullback of the original
cover. Combining this with \eqref{ds:permutation-invariants}, the
pullback square gives
\[
 \cL(V)^{hG}
 \simeq\cL(Y)\mathop{\times}_{\cL(Z)}\cL(Z)
 \simeq\cL(Y).
\]
Nilinvariance restores $X\times_S Y$ in place of $V$. In
dimension zero there are no proper nonempty coincidence loci, and
the same proof is simply the calculation for a disjoint union of
section copies. This starts the induction and proves
\eqref{ds:comparison}.

Taking $Y=S$ and using \eqref{ds:support-identification} yields
\eqref{ds:spectral-descent}. Finally, exact rational invariants give
\[
 \pi_0\Bigl(\bigl(\K^J(H\otimes_A^{\mathbf L}T)_\Q\bigr)^{hG}\Bigr)
   =\bigl(K_0^J(H\otimes_A^{\mathbf L}T)\otimes\Q\bigr)^G,
\]
which proves \eqref{ds:kzero-descent}, in fact as an isomorphism.
\end{proof}

\subsection{Comparison of ordinary and normalized Euler lengths}

We now isolate the consequence needed before introducing asymptotic
categories. Let $(A,\m_A,k)$ be a complete regular local
$\Z_{(p)}$-algebra of dimension $e>0$, with residue characteristic
$p$, and let $z_1,\ldots,z_e$ be regular parameters. Let $T$ be a
finite local domain extension of $A$ with the same residue field.
Put $r=\rank_A T$. Suppose that $C$ is a $T$-algebra equipped with
an additive nonnegative length $\lambda$ on $C$-modules annihilated
by powers of $\m_A$, with $\lambda(0)=0$. Infinite length is
allowed for general modules. Assume that, for every $v\geq1$,
\begin{equation}
 \lambda H_i(z_1^v,\ldots,z_e^v;C)=0\quad(i>0),
 \qquad
 \lambda\bigl(C/(z_1^v,\ldots,z_e^v)C\bigr)=v^e r.
 \label{ds:length-contract}
\end{equation}
Theorem~\ref{len:input} supplies exactly these hypotheses in our
applications, since the multiplicity of the base parameters on
$T$ is its generic rank over the regular base.
The next lemma is the formal version of
Corollary~\ref{len:finite-length-base-change}, using only
\eqref{ds:length-contract}.

\begin{lemma}
\label{ds:finite-length-basechange}
For every finite-length $A$-module $V$,
\begin{equation}
 \lambda\Tor_i^A(V,C)=0\quad(i>0),
 \qquad
 \lambda(V\otimes_A C)=r\length_A(V).
 \label{ds:finite-length-formula}
\end{equation}
All these lengths are finite.
\end{lemma}

\begin{proof}
The Koszul resolution on $z_1,\ldots,z_e$ resolves $k$ over
$A$, so \eqref{ds:length-contract} with $v=1$ proves the
assertion for $V=k$. Induct on the length of $V$ using a
composition series. In the long exact Tor sequence associated to
a short exact sequence with quotient $k$, nonnegativity and
additivity show that every positive Tor term has length zero.
The image of the connecting map into degree zero also has length
zero. Additivity in degree zero then gives the asserted factor
$r$. This simultaneously proves finiteness at every step.
\end{proof}

If a finite group $G$ acts on $T$ over $A$, then for $g\in G$ scalar
twisting means $g^*P=P\otimes_{T,g}T$,
where the left $T$-action on the last factor is through $g$.
On a free complex this applies $g$ to every differential matrix entry.

\begin{lemma}[Ordinary norm comparison]
\label{ds:ordinary-norm}
Assume, in addition, that a finite group $G$ acts on $T$ over $A$
and is transitive on every geometric fiber of $\Spec T\to\Spec A$.
For a perfect $T$-complex $P$ supported on $V(\m_A T)$, the
quantities
\[
 \chi_T(P)=\sum_i(-1)^i\length_T H_i(P),
 \qquad
 \chi_C(P)=\sum_i(-1)^i
           \lambda H_i(P\otimes_T^{\mathbf L}C)
\]
are finite and define homomorphisms on $K_0^{\m_A}(T)$.
They agree on every norm sum:
\begin{equation}
 \sum_{g\in G}\chi_T(g^*P)
   =\sum_{g\in G}\chi_C(g^*P).
 \label{ds:ordinary-norm-identity}
\end{equation}
More generally, the two homomorphisms agree on the image of
$K_0^{\m_A}(A)$ and on every $G$-invariant rational class.
\end{lemma}

\begin{proof}
We first establish finiteness and additivity, then compare the
evaluations on pullbacks before applying descent to the norm class.
Represent $P$ by a bounded finite free $T$-complex. Its ordinary
homology is finite and supported at the closed point, hence has
finite length. For the $C$-evaluation, choose a common $v$ such
that each $z_j^v$ acts nullhomotopically on $P$. This is possible
by the support criterion above. Set
$K_v=K(z_1^v,\ldots,z_e^v;T)$. Successively splitting the cones
of these nullhomotopic actions gives
\begin{equation}
 K_v\otimes_T P\otimes_T C
 \simeq
 \bigoplus_{j=0}^e
       (P\otimes_T C)[j]^{\oplus\binom ej}.
 \label{ds:koszul-splitting}
\end{equation}
The homotopies may be tensored with $C$ and with the other Koszul
factors, so this splitting respects the module structure on which
$\lambda$ is defined.

In the opposite order, taking Koszul homology first gives a finite
spectral sequence whose terms are finite direct sums of
$H_j(z_1^v,\ldots,z_e^v;C)$. Their lengths are finite by
\eqref{ds:length-contract}. Hence every homology module on the
left side of \eqref{ds:koszul-splitting} has finite length.
The splitting and nonnegativity then give finiteness for every
$H_i(P\otimes_T C)$. The homology modules and spectral-sequence terms
to which $\lambda$ is applied are annihilated by fixed powers of
$\m_A$. The Euler sums are therefore defined, and the long exact
homology sequence of a triangle proves their additivity.

Let $Q$ now be a perfect $A$-complex supported at $\m_A$.
Its homology modules have finite length. The finite hyper-Tor
spectral sequence for $Q\otimes_A^{\mathbf L}C$, together with
Lemma~\ref{ds:finite-length-basechange}, gives
\begin{equation}
 \chi_C(Q\otimes_A^{\mathbf L}T)
       =r\sum_j(-1)^j\length_A H_j(Q).
 \label{ds:normalized-base-euler}
\end{equation}
Indeed every positive Tor row has length zero and the zeroth row
has the displayed scaled lengths.

For the ordinary evaluation, choose a bounded finite free
$A$-resolution $L_\bullet$ of $T$, possible because $A$ is
regular. Its alternating rank is $r$, as is seen after tensoring
with the fraction field of $A$. For every finite-length
$A$-module $V$, the finite complex $V\otimes_A L_\bullet$
therefore gives
\[
 \sum_i(-1)^i\length_A\Tor_i^A(V,T)
   =\sum_i(-1)^i\rank_A(L_i)\length_A(V)
   =r\length_A(V).
\]
The hyper-Tor spectral sequence for
$Q\otimes_A^{\mathbf L}T$ consequently yields
\begin{equation}
 \chi_T(Q\otimes_A^{\mathbf L}T)
       =r\sum_j(-1)^j\length_A H_j(Q).
 \label{ds:ordinary-base-euler}
\end{equation}
We have used that finite-length $T$-modules have the same
$A$-length and $T$-length, because both residue fields are $k$.
Equations~\eqref{ds:normalized-base-euler} and
\eqref{ds:ordinary-base-euler} prove agreement on classes pulled
back from $A$.

Apply Theorem~\ref{ds:descent} with $H=A$ and $J=\m_A$.
Every invariant class in $K_0^{\m_A}(T)\otimes\Q$ is the
pullback of a rational class downstairs. Extending the two Euler
homomorphisms $\Q$-linearly, they therefore agree on that
invariant subspace. The class $\sum_{g\in G}[g^*P]$ is
invariant, which gives \eqref{ds:ordinary-norm-identity}.
\end{proof}

\begin{remark}
\label{ds:no-length-invariance}
No invariance of $\lambda$ under $G$ is asserted or needed.
The equality is for the sum of the conjugate classes, not for each
individual class. This distinction allows the two length algebras
used later to be chosen independently.
\end{remark}

\section{Supported asymptotic complexes}
\label{sec:asymptotic}

The goal is an analogue of Lemma~\ref{ds:ordinary-norm} for sequences of
bounded free complexes modulo free sequences of negligible normalized
rank, proved in Theorem~\ref{as:norm}. Two points require care: support
may arise only after passing to the quotient, and supported objects may
be idempotent summands whose projectors do not lift before quotienting.
Their homology lengths therefore cannot in general be computed from the
raw homologies of arbitrary lifts. We first construct nonnegative lengths
on uniformly supported lifts, extend them through the supported quotient
and its idempotent completion, and prove their compatibility with scalar
extension on every supported summand.

Throughout this section, $(A,J)$ is an excellent regular local
$\Z_{(p)}$-algebra of dimension $e>0$ with residue characteristic $p$,
where $J$ is the maximal ideal, and $z_1,\ldots,z_e$ is a regular system
of parameters. Let $T$ be a finite
local domain over $A$ with the same residue field, and put
$r=\rank_A T$. Let $C$ be a $T$-algebra equipped with a nonnegative
additive length $\lambda$ on its modules annihilated by powers of $J$.
The assumptions on this length that are used in this section are
\begin{equation}
 \lambda H_j(z_1^v,\ldots,z_e^v;C)=0\quad(j>0),
 \qquad
 \lambda\bigl(C/(z_1^v,\ldots,z_e^v)C\bigr)=v^e r
 \quad(v\geq 1).
 \label{as:length-hypotheses}
\end{equation}
The length is allowed to vanish on nonzero modules. In the applications,
Theorem~\ref{len:input} supplies these hypotheses, as well as the stronger
parameter properties used in the other sections. Write $K_v$ for the
Koszul complex $K(z_1^v,\ldots,z_e^v;A)$, in homological degrees
$0,\ldots,e$.

\subsection{The rank quotient and its support condition}

Fix positive integers $s_n$ and a nonprincipal ultrafilter $\mathcal U$
on the positive integers. Rank estimates $O(s_n)$ and $o(s_n)$ always
refer to the ordinary limit, not to $\mathcal U$. If $(a_n/s_n)$ is
bounded, write
\[
 \overline a=\lim_{n\to\mathcal U}\frac{a_n}{s_n}.
\]
For $S=A$ or $T$, consider the additive category of sequences
$(S^{a_n})_n$ with $a_n=O(s_n)$, whose morphisms are arbitrary
sequences of matrices over $S$. Let $\mathcal D_S$ be its perfect
hull: take complexes in one common bounded degree range for the whole
sequence, and then idempotent completion. Let
$\mathcal N_S$ be the thick subcategory generated by the free
sequences with $a_n=o(s_n)$, and define
\[
 \mathcal P_S=(\mathcal D_S/\mathcal N_S)^{\natural}.
\]
Here $(-)^{\natural}$ denotes idempotent completion. The quotient
functor is denoted by $q_S$.

A thick subcategory is closed under shifts, cones, and direct summands.
The Verdier quotient inverts maps whose cones lie in that subcategory;
idempotent completion then adjoins the images of idempotent endomorphisms.

An object $X$ in $\mathcal D_S$ or $\mathcal P_S$ is called
\emph{supported on $J$} if there is an integer $v\geq1$ such that
$z_j^v\operatorname{id}_X=0$ for every $j$. The identity is taken
in the category containing $X$. Thus in $\mathcal D_S$ this is a
uniform homotopical support condition on the sequence, whereas in
$\mathcal P_S$ the actions need vanish only after the rank-small
quotient. Write $\mathcal S_S\subseteq\mathcal D_S$ and
$\mathcal P_S^J\subseteq\mathcal P_S$ for the resulting full
subcategories. They are thick: nilpotence of central actions is
preserved by shifts, cones, and retracts, allowing larger powers
for cones. We identify this condition with central-localization
support before applying descent at the end of the section.

\begin{remark}
\label{as:nonuniform-warning}
Raw homology lengths need not survive the rank quotient. Over a
discrete valuation ring $A=T$ with parameter $z$, take $s_n=n$ and
the rank-one free resolution of $T/(z^n)$. This sequence belongs to
$\mathcal N_T$, whereas its raw normalized degree-zero homology
length is $1$. It has no uniform annihilating power before
quotienting. For each fixed $v$, however, its homology lengths after
tensoring with $K(z^v;T)$ are bounded independently of $n$ and hence
become zero after normalization. This is why we construct homology
lengths on uniformly supported lifts first, and then extend them to
supported quotient objects.
\end{remark}

\subsection{The supported subquotient}

The next lemma supplies that extension's categorical starting point:
every supported quotient object is a summand of a uniformly supported
lift, and morphisms between such lifts can be computed using supported
roofs. Neither assertion requires the summand projector to lift.

\begin{lemma}
\label{as:supported-quotient}
The natural exact functor
\[
 \mathcal S_S/(\mathcal S_S\cap\mathcal N_S)
 \longrightarrow \mathcal D_S/\mathcal N_S
\]
is fully faithful. Its idempotent completion is
$\mathcal P_S^J$.
\end{lemma}

\begin{proof}
First let $X\in\mathcal S_S$, $Y\in\mathcal N_S$, and
$f:X\to Y$. Choose $v$ such that $z_j^v\operatorname{id}_X=0$
for every $j$. Centrality gives $z_j^vf=0$. The exact Hom sequence
of the triangle
\[
 \operatorname{fib}(z_j^v:Y\to Y)\longrightarrow Y
 \xrightarrow{z_j^v}Y
\]
therefore lifts $f$ through the first object. Repeat this operation
for all the parameters, always using the zero action on the original
source $X$. The result factors $f$ through
$K_v\otimes_A Y[-e]$. This object belongs to $\mathcal N_S$ by
thickness, and to $\mathcal S_S$ by the usual Koszul homotopies.
Thus every map from a supported object to a small object factors
through a supported small object.

Here is the resulting roof argument. A morphism between $X,Y$ in
$\mathcal S_S$, computed in $\mathcal D_S/\mathcal N_S$, has a
representative
\[
 X\xleftarrow{s}Z\xrightarrow{f}Y,
 \qquad \operatorname{cofib}(s)=N_0\in\mathcal N_S.
\]
In the triangle $Z\to X\xrightarrow{u}N_0$, factor $u$ through
$N_1\in\mathcal S_S\cap\mathcal N_S$ by the preceding paragraph.
Let $Z_1$ be the fiber of $X\to N_1$. The induced morphism of
triangles gives $c:Z_1\to Z$ and $s_1:Z_1\to X$ with $sc=s_1$.
Now $Z_1\in\mathcal S_S$ and $\operatorname{cofib}(s_1)=N_1$,
so $X\xleftarrow{s_1}Z_1\xrightarrow{fc}Y$ is a representative
entirely in the supported subquotient. This proves fullness.

For faithfulness, suppose a map $f:Z\to Y$ between supported
objects becomes zero in $\mathcal D_S/\mathcal N_S$. The calculus
of Verdier roofs gives a map $t:W\to Z$ with small cone and $ft=0$.
The triangle of $t$ makes $f$ factor through that small cone.
Factoring its first leg through a supported small object as above
shows that $f$ is already zero in the supported subquotient. Apply
this to the numerator of any roof to obtain faithfulness.

Finally take $X\in\mathcal P_S^J$. It is a retract of $q_S(Y)$
for some $Y\in\mathcal D_S$. Choose $v$ with all $z_j^v$ acting
as zero on $X$. The cone of each such zero map splits, so $X$ is a
retract of $K_v\otimes_A X$. The latter is a retract of
$q_S(K_v\otimes_A Y)$, whose lift is uniformly supported. Hence
$X$ belongs to the idempotent completion of the fully faithful
supported subquotient. The reverse inclusion follows because
support is closed under retracts.
\end{proof}

\subsection{Nonnegative homology lengths}

Every supported quotient object is now accessible through a uniformly
supported lift and a summand. We next construct its nonnegative homology
lengths and show that they give additive Euler evaluations.

For ordinary length, let $\mathcal A_T$ be the abelian category
of sequences of finite $T$-modules annihilated by some common power
of $J$ and with lengths $O(s_n)$. For the $C$ evaluation, let
$\mathcal A_C$ be the category of sequences of $C$-modules with
the same uniform torsion condition and with finite $\lambda$-lengths
$O(s_n)$. The power and the implied constant can depend on the
object. Kernels, cokernels, and extensions preserve these conditions,
so both categories are abelian. In each case the sequences of
normalized ultralimit length zero form a Serre subcategory. Denote
the Serre quotients by $\overline{\mathcal A}_T$ and
$\overline{\mathcal A}_C$, and their nonnegative additive lengths
by $\overline\length_T$ and $\overline\lambda$, respectively.
These lengths are finite on every object.
Here negligible length means zero ultralimit; the free sequences
generating $\mathcal N_T$ still satisfy the stronger ordinary estimate
$o(s_n)$. We will show that supported objects killed by the rank
quotient have zero homology in these abelian quotients.

\begin{proposition}
\label{as:evaluations}
There are homological functors
\[
 \mathsf H_i^{\mathrm{ord}}:\mathcal P_T^J
       \longrightarrow\overline{\mathcal A}_T,
 \qquad
 \mathsf H_i^C:\mathcal P_T^J
       \longrightarrow\overline{\mathcal A}_C
\]
whose values on a uniformly supported lift are its ordinary homology
sequence and its homology sequence after tensoring with $C$,
respectively. For each object $X$, their lengths
\[
 b_i^{\mathrm{ord}}(X)=\overline\length_T
       \mathsf H_i^{\mathrm{ord}}(X),
 \qquad
 b_i^C(X)=\overline\lambda\mathsf H_i^C(X)
\]
are finite nonnegative real numbers, zero outside a finite range.
Their Euler sums define homomorphisms
\[
 \chi_{\mathrm{ord}},\chi_C:K_0(\mathcal P_T^J)\longrightarrow\mathbb R.
\]
If $F=(F_n)$ is a bounded free sequence whose image is supported,
and $v$ makes all $z_j^v$ act as zero on that image, then for either
choice of evaluation, writing $b_i(F)=b_i(q_TF)$,
\begin{equation}
 (1+t)^e\sum_i b_i(F)t^i
  =\sum_i t^i\lim_{n\to\mathcal U}
     \frac{\operatorname{len}
       H_i(K_v\otimes_A F_n)}{s_n}.
 \label{as:koszul-polynomial}
\end{equation}
Here $\operatorname{len}=\length_T$ in the ordinary case; in the
$C$ case the complex on the right is further tensored over $T$
with $C$, and $\operatorname{len}=\lambda$. The identity is one
of finite Laurent polynomials.
\end{proposition}

\begin{proof}
We first bound homology on $\mathcal S_T$, then show that it kills
$\mathcal S_T\cap\mathcal N_T$, and finally extend it to the
idempotent-completed quotient.

Begin with a bounded free sequence $F$ in $\mathcal S_T$. Choose
uniform powers annihilating its identity. Splitting their cones
gives an isomorphism in the homotopy category
\begin{equation}
 K_v\otimes_A F\simeq
 \bigoplus_{j=0}^e F[j]^{\oplus\binom ej}.
 \label{as:split-koszul}
\end{equation}
For ordinary length, every Koszul homology $H_j(K_v\otimes_A T)$
has finite length. Taking Koszul homology first in the bounded
double complex bounds the homology lengths of $K_v\otimes_A F_n$
by fixed constants times the ranks of the terms of $F_n$. They
are therefore $O(s_n)$. For $C$, the same spectral sequence uses
Equation~\eqref{as:length-hypotheses}: positive Koszul rows have
length zero and the zero row has finite length. Additivity and
nonnegativity give the same bound. The splitting in
Equation~\eqref{as:split-koszul} bounds the homologies of $F_n$
themselves. They have a common primary annihilator, because a power
of every $z_j$ annihilates them.

A general object $X\in\mathcal S_T$ may be a retract in
$\mathcal D_T$ of a bounded free sequence $F$ that is not supported.
Choose $v$ killing the parameter actions on $X$. Then $X$ is a retract
of $K_v\otimes_A X$, which is a retract of the uniformly supported
bounded free sequence $K_v\otimes_A F$. The preceding bounds apply
to this sequence and hence to the homology summands defining $X$.
This is a retract before taking the rank quotient. Thus homology on
$\mathcal S_T$ takes
values in the indicated abelian length quotients.

Fix $v$. Koszul tensoring a rank-small free generator has homology
of normalized length zero by the same bounds. The class of objects
$Y\in\mathcal D_T$ with this property for $K_v\otimes_A Y$
is thick: use homology long exact sequences, nonnegativity, and
retracts. It therefore contains $\mathcal N_T$. If
$X\in\mathcal S_T\cap\mathcal N_T$, choose $v$ killing its
parameter actions. Then $X$ is a retract of $K_v\otimes_A X$,
so its homology is zero in both length quotients.

The homology functors consequently invert maps whose cones are
supported and small. They descend to
$\mathcal S_T/(\mathcal S_T\cap\mathcal N_T)$, preserving
homology long exact sequences. Now consider summands after quotienting.
By Lemma~\ref{as:supported-quotient},
every object of $\mathcal P_T^J$ is a summand of an object of this
subcategory. Extend a homology functor to a summand $(X,a)$ by
taking $\operatorname{im}\mathsf H_i(a)$ in its abelian target.
This does not require a sequence of chain projectors lifting $a$.
Distinguished triangles in an idempotent completion are summands
of distinguished triangles in the original triangulated category
\cite[Definition~1.10 and Theorem~1.12]{BalmerSchlichting}.
Taking images of compatible idempotents preserves exact sequences
in an abelian category, so the extended functors are still
homological.

Each object is a summand of a fixed bounded uniformly supported
Koszul lift. Its homology objects therefore vanish outside a finite
range, and their lengths are bounded by those of that lift. This
proves finiteness, nonnegativity, and boundedness of the $b_i$.
The finite long exact sequences prove additivity of the Euler sums.

For the last assertion, the splitting in
Equation~\eqref{as:split-koszul} now takes place in $\mathcal P_T$,
since the chosen powers act as zero there. Applying the extended
homological functors gives the left side of
Equation~\eqref{as:koszul-polynomial}. Its Koszul tensor has a
uniformly supported lift, whose actual homology computes the right
side. This proves the identity without defining its coefficients
by formal division.
\end{proof}

In particular, Equation~\eqref{as:koszul-polynomial} assigns zero
to the quotient object in Remark~\ref{as:nonuniform-warning}, as
required. Its nonnegative coefficients come from homological
functors, not from an assumed positivity of a polynomial quotient.

\subsection{Scalar extension and nonlifting idempotents}

Construct $\overline{\mathcal A}_A$ and its length in the same way
using ordinary finite $A$-modules. The preceding construction over
$A$ with ordinary length defines $\chi_A$ on $\mathcal P_A^J$.
Scalar extension gives an exact functor
$\mathcal P_A^J\to\mathcal P_T^J$.

The supported Euler evaluations are defined on every object of
$\mathcal P_T^J$. To apply descent, it remains to prove that they agree
on all classes coming from $\mathcal P_A^J$, including idempotent
summands.

The normalized evaluation uses the following consequence of the
length hypotheses. We state it for the present base $A$, which need
not be complete.

\begin{lemma}
\label{as:finite-tor}
For every finite-length $A$-module $V$,
\[
 \lambda\Tor_j^A(V,C)=0\quad(j>0),
 \qquad
 \lambda(V\otimes_A C)=r\length_A V.
\]
All these modules are annihilated by a power of $J$.
\end{lemma}

\begin{proof}
The Koszul and composition-series proof of
Lemma~\ref{ds:finite-length-basechange} applies here: it uses regularity
and the length hypothesis, but not completeness. A power of $J$
annihilating $V$ also annihilates its derived tensor homology.
\end{proof}

\begin{proposition}
\label{as:base-equality}
For every $X\in\mathcal P_A^J$,
\[
 \chi_{\mathrm{ord}}(X\otimes_A^{\mathbf L}T)
   =r\chi_A(X)
   =\chi_C(X\otimes_A^{\mathbf L}T).
\]
These equalities hold for summands defined by idempotents in the
Verdier quotient, without assuming that the idempotents lift to
chain maps before localization.
\end{proposition}

\begin{proof}
The two evaluations recover the factor $r$ differently: the ordinary
one uses the alternating rank of a resolution of $T$, while the
normalized one uses the vanishing of positive Tor lengths. To make
both calculations on an arbitrary summand, we descend their finite
spectral sequences through supported Verdier roofs before taking
idempotent images.

Let $Y$ be a uniformly supported bounded free sequence over $A$ and
let $a$ be an idempotent endomorphism of its image in
$\mathcal S_A/(\mathcal S_A\cap\mathcal N_A)$. Every supported
object of $\mathcal P_A$ has this form as a summand, by
Lemma~\ref{as:supported-quotient}: in its proof one may enlarge the
initial lift to a bounded free sequence before taking the Koszul
tensor. Put
\[
 V_b=\operatorname{im}\mathsf H_b^A(a)
       \quad\text{in }\overline{\mathcal A}_A.
\]
These images exist in the abelian quotient even when $a$ has no
chain-level projector representing it.

Choose a fixed bounded finite free $A$-resolution $P_\bullet$ of
$T$; such a resolution exists because $A$ is regular. The double
complex $Y\otimes_A P_\bullet$ computes
$Y\otimes_A^{\mathbf L}T$. Filtering by the degree of $P$ gives
a finite spectral sequence with first page
\begin{equation}
 E^1_{u,b}=H_b(Y)\otimes_A P_u.
 \label{as:ordinary-page}
\end{equation}
It is natural in homotopy-category maps from this page onward:
homotopies in $Y$ preserve the filtration and act trivially on
the first page. On retracts in the prequotient category the same
construction is defined by taking the corresponding idempotent
images on every page and on the abutment filtration. Its terms
and finite abutment filtration can be
regarded as $A$-modules. The abutment's ordinary $A$-length equals
its $T$-length, since the residue fields agree.

A map with supported small cone is an isomorphism on $H_b(Y)$
in $\overline{\mathcal A}_A$. It is thus an isomorphism on
Equation~\eqref{as:ordinary-page}, on every subsequent page, and
on the finite filtered abutment in the length quotient. The
spectral sequence therefore descends through all Verdier roofs
in the supported subquotient. The natural augmentation
$\operatorname{Tot}(Y\otimes_A P_\bullet)\to Y\otimes_A T$
identifies its abutment with the homology of scalar extension,
and this identification also descends through those roofs.
Thus the induced action of $a$ on the abutment is precisely the
scalar-extended idempotent, not a separately chosen action.

In particular $a$ acts as an actual
idempotent on every page and on that filtered abutment. The
differentials commute with this action. Taking its images commutes
with passage to homology, since an idempotent decomposes a complex
in an abelian category into its image and complementary image.
It also preserves the finite exact sequences defining the
abutment filtration: for a filtration preserved by an idempotent
$a$, one has $\operatorname{im}(a)\cap F_u=a(F_u)$ and hence
$\operatorname{gr}_u\operatorname{im}(a)
=\operatorname{im}(\operatorname{gr}_u a)$.
Consequently the Euler length on the
scalar-extended summand is the Euler length of the images on
the first page. They are $\rank_A(P_u)$ copies of $V_b$, and hence
this value is
\begin{equation}
 \left(\sum_u(-1)^u\rank_A P_u\right)
 \sum_b(-1)^b\overline\length_A V_b
 =r\sum_b(-1)^b\overline\length_A V_b.
 \label{as:ordinary-idempotent}
\end{equation}
The equality of alternating ranks with $r$ follows by tensoring
the resolution with the fraction field of $A$.

For $C$, use the finite hyper-Tor spectral sequence
\begin{equation}
 {}^CE^2_{u,b}=\Tor_u^A(H_b(Y),C)
 \ \Longrightarrow\ H_{u+b}(Y\otimes_A^{\mathbf L}C).
 \label{as:c-page}
\end{equation}
It can be constructed from the functorial finite Postnikov
filtration of $Y$. It has bounded $b$ and $0\leq u\leq e$, so
there is no unbounded convergence issue. Lemma~\ref{as:finite-tor}
makes every positive Tor term zero in
$\overline{\mathcal A}_C$. It also implies that
$V\mapsto V\otimes_A C$ induces an exact functor
\[
 \overline{\mathcal A}_A\longrightarrow
 \overline{\mathcal A}_C
\]
that multiplies length by $r$. Indeed, on a short exact sequence
its possible failure of left exactness is the image of a
positive Tor module of zero length. Negligible-length sequences
map to negligible-length sequences by the same lemma. This
explains both exactness after the Serre quotient and descent of
the functor to that quotient.

It follows that a supported-small-cone map induces isomorphisms
on the second page of Equation~\eqref{as:c-page} and its finite
filtered abutment in the length quotient. Thus this spectral
sequence too is functorial through the supported Verdier roofs.
Taking the images of the same idempotent $a$ gives only the zero
Tor row, namely the images of $V_b$ under this exact tensor
functor. Their Euler length is
$r\sum_b(-1)^b\overline\length_A V_b$, exactly the value in
Equation~\eqref{as:ordinary-idempotent}. Finally this last sum
is $\chi_A(Y,a)$ by construction. This proves both equalities
for every summand and hence for every $X\in\mathcal P_A^J$.
\end{proof}

\subsection{Connective coefficients and the norm comparison}

We have constructed the supported Euler evaluations and proved their
agreement after scalar extension from $A$. To extend that agreement
to conjugate sums, we now put the rank quotient in the form required
by Theorem~\ref{ds:descent}. The coefficient algebra must be
connective, and the support condition must be central-localization
support.

\begin{lemma}[Connective coefficients]
\label{as:connective}
There is a connective associative derived $A$-algebra $H$ and
compatible equivalences
\[
 \mathcal P_A\simeq\Perf(H),
 \qquad
 \mathcal P_T\simeq\Perf(H\otimes_A^{\mathbf L}T).
\]
The equivalences respect central localization by elements of $A$.
\end{lemma}

\begin{proof}
The sequence $E_S=(S^{s_n})_n$ is an additive generator. Indeed, if
$a_n=O(s_n)$, there is a single integer $c$ such that $a_n\leq cs_n$
for every $n$, after increasing $c$ to handle the finitely many
exceptional indices. Coordinate inclusions and projections make
$(S^{a_n})_n$ a summand of $E_S^c$. Consequently
\[
 \mathcal D_S\simeq\Perf(B_S),
 \qquad B_S=\prod_n M_{s_n}(S).
\]

The $A$-module underlying $B_A$ is flat. One can check this by the
ideal criterion: every finitely generated ideal of the Noetherian
ring $A$ is finitely presented, so its tensor product with this
product of finite free modules is the corresponding product of its
tensor products; each required injection is then a product of
injections. Moreover, $T$ is finitely presented over $A$, so tensoring
with $T$ commutes with that product. Thus
\begin{equation}
 B_A\otimes_A^{\mathbf L}T
 \simeq B_A\otimes_A T
 \cong B_T.
 \label{as:product-base-change}
\end{equation}

The same product argument makes every morphism module between the
original free sequences flat over $A$. Hence the explicit dg quotient
of this degree-zero category by its full rank-small subcategory computes
the derived quotient \cite[Section~3.1 and Theorem~3.4]{Drinfeld}.
Its morphisms are sums of paths formed from tensor products of the
original morphism modules and contracting homotopies in homological
degree $1$. They are concentrated in nonnegative degrees and degreewise
$A$-flat, hence are $K$-flat. Passing to the perfect hull and idempotent
completion kills exactly $\mathcal N_A$. The image of $E_A$ remains a
generator, and its endomorphism algebra
$H=\operatorname{End}_{\mathcal P_A}(q_AE_A)$ is connective, giving
$\mathcal P_A\simeq\Perf(H)$. This assertion concerns the chosen
generator, not arbitrary shifts of perfect objects.

Since $T$ is finitely presented over $A$, tensoring each original
morphism module with $T$ gives the corresponding module of morphisms
between free $T$-sequences. Tensoring the path formula therefore gives
exactly the $T$-linear dg quotient by the rank-small free sequences.
The same flatness criterion applies over the Noetherian ring $T$.
The upstairs endomorphism algebra is consequently
$H\otimes_A T\simeq H\otimes_A^{\mathbf L}T$, and its perfect hull
and idempotent completion are $\mathcal P_T$. This also agrees with
the preservation of exact quotients under scalar extension
\cite[Lemma~3.12]{LT}. All constructions are centrally $A$-linear and
preserve central localization.
\end{proof}

Under these equivalences, the support condition used throughout this
section is exactly the one in Theorem~\ref{ds:descent}. Indeed, for
a perfect module $X$, localization at a central element $z_j$ is
the telescope of multiplication by $z_j$ in the module category.
Compactness implies that $X[1/z_j]=0$ if and only if some power
$z_j^v$ kills its identity. A common power can be chosen for the
finitely many parameters. The same argument for $\Perf(B_S)$
identifies $\mathcal S_S$ with the supported objects before the
rank quotient. In particular, this comparison does not require a
uniformly supported lift of an object in $\mathcal P_S^J$.

Suppose now that a finite group $G$ acts on $T$ over $A$, that
$\Spec T\to\Spec A$ is surjective, and that $G$ is transitive
on every geometric fiber, as in Theorem~\ref{ds:descent}. It acts
on $\mathcal P_T^J$ by scalar twisting. We write $gX$ for the
functor induced by $g^*P=P\otimes_{T,g}T$, with the convention of
Section~\ref{sec:descent}.

\begin{theorem}[Asymptotic norm comparison]
\label{as:norm}
Let $A,J,T,C,\lambda$ satisfy the setup of this section and
\eqref{as:length-hypotheses}. Suppose a finite group $G$ acts on $T$
over $A$, the map $\Spec T\to\Spec A$ is surjective, and $G$ is
transitive on every geometric fiber. Then, for every
$X\in\mathcal P_T^J$,
\[
 \sum_{g\in G}\chi_{\mathrm{ord}}(gX)
   =\sum_{g\in G}\chi_C(gX).
\]
Neither invariance of $\lambda$ under $G$ nor equality of the two
evaluations on individual classes is assumed.
\end{theorem}

\begin{proof}
Lemma~\ref{as:connective} identifies the two supported categories
with the supported perfect categories of $H$ and
$H\otimes_A^{\mathbf L}T$, where $H$ is connective.
The identification over $T$ is compatible with the $G$-action:
scalar twisting acts on the $T$ factor in the coefficient algebra.
Theorem~\ref{ds:descent} therefore gives a surjection
\[
 K_0(\mathcal P_A^J)\otimes\Q
 \longrightarrow
 \bigl(K_0(\mathcal P_T^J)\otimes\Q\bigr)^G.
\]
The norm class $\sum_g[gX]$ is invariant. The homomorphism
$\chi_{\mathrm{ord}}-\chi_C$, extended rationally to the
Grothendieck group, vanishes on the image from $A$ by
Proposition~\ref{as:base-equality}. Applying it to a preimage
of this norm class proves the assertion.
\end{proof}

\section{Proof of strict positivity}\label{sec:positivity}

We apply the two norm comparisons to different finite extensions.  The
first comparison replaces one of the original modules by a normalized
length algebra.  The second comparison is applied to bounded complexes
obtained from finite approximations to a parameter quotient of that
algebra.  The normalized lengths need not be invariant under either
Galois group.

\subsection{Reduction and the first norm comparison}

\begin{lemma}\label{pos:reduction}
To prove Theorem~\ref{thm:main}, it suffices to prove the assertion for
\[
 D=R/P,\qquad E=R/Q,\qquad
 h=\dim D>0,\qquad l=\dim E>0,\qquad h+l=d,
\]
where $R$ is complete with algebraically closed residue field $k$ of
characteristic $p>0$, $P,Q$ are prime, and $P+Q$ is primary to $\m_R$.
\end{lemma}

\begin{proof}
The equicharacteristic case is Serre's positivity theorem
\cite{Serre}.  We may therefore assume that the residue characteristic
is positive.  Completion, followed by the extension in
Lemma~\ref{len:residue-extension}, reduces to a complete regular local
ring with algebraically closed residue field.  For clarity, a flat
local extension $R\to R_1$ with $\m_RR_1=\m_{R_1}$ preserves the
length of every finite-length module: a composition series becomes a
filtration with factors $R_1/\m_{R_1}$.  Flat base change identifies the
Tor modules after extension.  The same argument applied to
$M/\m_R^nM$ preserves the Hilbert--Samuel function, and hence the
dimension of $M$; it applies equally to $N$.  It also preserves
regularity, since it preserves the dimension and the minimal number
of generators of the maximal ideal.

We use the vanishing theorem of Roberts and Gillet--Soul\'e in its
general regular local form \cite{Roberts,GilletSoule}:
\begin{equation}\label{pos:vanishing}
 \chi^R(V,W)=0
 \quad\text{if}\quad
 \length_R(V\otimes_R W)<\infty
 \quad\text{and}\quad \dim V+\dim W<d.
\end{equation}
Take prime filtrations of $M$ and $N$.  All pairs of factors still
have finite-length tensor product, since their supports are contained
in the respective original supports.  Additivity of the Euler
characteristic and Equation~\eqref{pos:vanishing} discard every pair except
those having both original dimensions.  There are top-dimensional
factors in each filtration, so at least one such pair remains.
Proving strict positivity for each remaining pair proves it for
$M,N$.  Their primes satisfy $\sqrt{P+Q}=\m_R$.

Finally, a zero-dimensional prime quotient is $k$, and its
complementary $d$-dimensional prime quotient is $R$, since a regular
local ring is a domain.  For that pair the characteristic equals
$\length_R(k)=1$.  Thus we may assume $h,l>0$.
\end{proof}

Fix a pair as in Lemma~\ref{pos:reduction}.  By
Lemma~\ref{len:normalization}, choose regular complete normalization
bases
\[
 A_D\subset D,\qquad A_E\subset E
\]
of dimensions $h,l$, with residue field $k$, such that the extensions
are finite and generically separable.

Let $D'$ and $E'$ be the integral closures of these bases in finite
Galois closures of the respective fraction-field extensions.  They
are finite over their bases by excellence, and are finite over $D,E$
as well.  They are local: a finite algebra over a complete local ring
is a product of complete local algebras, and a domain has only one
factor \cite[Tag~03QH]{Stacks}. Their residue fields are finite extensions of $k$, hence are
$k$.  Write
\[
 G_D=\operatorname{Gal}(\operatorname{Frac}(D')/
                  \operatorname{Frac}(A_D)),\qquad
 G_E=\operatorname{Gal}(\operatorname{Frac}(E')/
                  \operatorname{Frac}(A_E)).
\]
Normality of the bases gives $(D')^{G_D}=A_D$ and
$(E')^{G_E}=A_E$. For either extension write $T^G=A$.
Let $\Omega$ be an algebraically closed $A$-field and put
$S=T\otimes_A\Omega$. For $u=\sum_i t_i\otimes\omega_i\in S$,
the norm $\prod_{g\in G}g(u)$ is scalar: the coefficients of
$\prod_{g\in G}(\sum_i g(t_i)X_i)$ lie in $T^G=A$, and we
specialize $X_i=\omega_i$. If the finite Artinian $\Omega$-algebra
$S$ had more than one orbit of points, its product decomposition
would give a nontrivial $G$-invariant idempotent $e$, equal to one on
one orbit and zero on the others. But its norm is $e^{|G|}=e$ and is
scalar, so $e=0$ or $e=1$, a contradiction. Thus every geometric fiber
is transitive, without any assertion that invariants commute with
base change. The prime and residue-field formulation is also given
in \cite[Tags~0BRI and 0BRJ]{Stacks}. Thus the hypotheses of
Theorem~\ref{ds:descent} apply to both extensions.

Use Theorem~\ref{len:input} to choose the algebras and normalized
lengths
\[
 (C_D,\lambda_D)=(C_{D'},\lambda_{D'}),\qquad
 (C_E,\lambda_E)=(C_{E'},\lambda_{E'}).
\]
They remain fixed throughout the proof.  Set
$a_D=\rank_D D'$ and $a_E=\rank_E E'$.  For $\sigma\in G_D$,
define
\[
 \iota_\sigma:R\longrightarrow D'
      \xrightarrow{\sigma}D'.
\]
Give $D'$ this $R$-action, and give $C_D$ the action obtained by
composing $\iota_\sigma$ with $D'\to C_D$. Define
\begin{equation}\label{pos:chi-sigma}
 c_\sigma
   =\sum_{i=0}^d(-1)^i
      \lambda_D\bigl(\Tor_i^R(C_D,E')\bigr),
\end{equation}
where this $R$-action on $C_D$ is understood.

\begin{lemma}\label{pos:first-average}
The lengths in Equation~\eqref{pos:chi-sigma} are finite, and
\begin{equation}\label{pos:first-average-equation}
 |G_D|a_Da_E\,\chi^R(D,E)=\sum_{\sigma\in G_D}c_\sigma.
\end{equation}
\end{lemma}

\begin{proof}
Take a finite $R$-free resolution of $E'$ and extend it to $D'$.
The resulting bounded free complex is supported at the closed point:
after localizing its homology away from $\m_R$, the supports of
$D'$ and $E'$ are disjoint.  The same holds after every twist, whose
kernel on $R$ is still $P$.  Lemma~\ref{ds:ordinary-norm} supplies
finite normalized homology lengths and equality of the ordinary and
normalized Euler evaluations after summing all conjugates.

Applying $\sigma$ entrywise gives a semilinear isomorphism from the
untwisted complex to its conjugate, so all ordinary Euler
characteristics in that sum are equal to $\chi^R(D',E')$.  Ordinary
lengths over $D'$ and over $R$ agree here because both residue fields
are $k$.

Choose an injection $D^{a_D}\to D'$ with torsion cokernel.  The
cokernel has dimension less than $h$, and therefore its intersection
characteristic with $E'$ vanishes by Equation~\eqref{pos:vanishing}.
Additivity gives
\[
 \chi^R(D',E')=a_D\chi^R(D,E').
\]
The analogous argument for $E\subset E'$ gives
$\chi^R(D,E')=a_E\chi^R(D,E)$.  The norm comparison now yields
Equation~\eqref{pos:first-average-equation}.
\end{proof}

It remains to show $c_\sigma>0$ for each $\sigma$.  Fix one
$\sigma$ until the final paragraph, and write $C=C_D$ and
$\lambda=\lambda_D$ with the chosen $R$-action.

\subsection{A parameter quotient and its Tor groups}

Choose
\begin{equation}\label{pos:parameters}
 x_1,\ldots,x_h\in Q
\end{equation}
which form an $R$-regular sequence and whose images are a system of
parameters of $D$.  Here is the simultaneous prime-avoidance
argument.  After $i-1$ choices, the associated primes of
$R/(x_1,\ldots,x_{i-1})$ have height $i-1$, since $R$ is
Cohen--Macaulay.  None contains $Q$, whose height is $h\ge i$.
Also avoid the primes in $D/(x_1,\ldots,x_{i-1})$ which obstruct the
next parameter dimension drop.  Each has positive dimension; if its
inverse image contained $Q$, it would contain $P+Q$ and hence would
be the maximal ideal, a contradiction.  Prime avoidance inside $Q$
therefore makes both choices possible at every step.

Put
\[
 B=R/(x_1,\ldots,x_h),\qquad
 I_\sigma=(\iota_\sigma(x_1),\ldots,
                 \iota_\sigma(x_h))D',\qquad U=C/I_\sigma C.
\]
The ideal $I_\sigma$ is a parameter ideal of $D'$, since finiteness
and then the automorphism $\sigma$ preserve that property.  Thus
Theorem~\ref{len:input} gives
\begin{equation}\label{pos:U-positive}
 0<\lambda(U)=e(I_\sigma,D')<\infty,
 \qquad
 \lambda\bigl(H_q(x;C)\bigr)=0\quad(q>0).
\end{equation}
Here the notation $x$ in a complex over $C$ uses $\iota_\sigma$.
Some integer $a\ge1$ satisfies
\begin{equation}\label{pos:uniform-annihilator}
 \m_B^a(D'/I_\sigma)=0.
\end{equation}
In particular it annihilates $U$ and every module over $D'/I_\sigma$.

\begin{lemma}\label{pos:tor-bound}
For every finite $B$-module $L$, all the following lengths are finite,
and
\begin{equation}\label{pos:tor-identification}
 \tau_i(L):=\lambda\bigl(\Tor_i^B(U,L)\bigr)
     =\lambda\bigl(\Tor_i^R(C,L)\bigr),
 \qquad \tau_i(L)=0\quad(i>l).
\end{equation}
\end{lemma}

\begin{proof}
The complexes in question have compatible $C$-actions.  The Koszul
homology $H_q(x;C)$ is annihilated by $x$ and by $P$, and is therefore
bounded-primary-torsion.  For $q>0$ its length is zero by
Equation~\eqref{pos:U-positive}.  A free $B$-resolution of $L$,
with each term of finite rank, gives
\[
 \lambda\bigl(\Tor_p^B(H_q(x;C),L)\bigr)=0\quad(q>0).
\]
Indeed every such homology is a subquotient of a finite direct sum
of a length-zero module.  For $q=0$ these lengths are finite because
$\lambda(U)<\infty$.  The change-of-rings spectral sequence
\[
 \Tor_p^B\bigl(\Tor_q^R(C,B),L\bigr)
       \ \Longrightarrow\ \Tor_{p+q}^R(C,L)
\]
now gives the asserted equality of lengths.  Only finitely many
terms enter each total degree.

For the sharper bound, choose an $R$-free resolution $P_L$ of $L$
in degrees $0,\ldots,d$ and consider
\[
 K(x;R)\otimes_R C\otimes_R P_L.
\]
Taking Koszul homology first, all positive rows have length zero,
and the zero row is $U\otimes_R P_L$.  Consequently the total
homology has length zero in degrees greater than $d$.  On the other
hand, $xL=0$, so in the derived category of $R$-modules
\[
 K(x;R)\otimes_R^{\mathbf L}L
       \simeq\bigoplus_{j=0}^h L[j]^{\binom hj}.
\]
Tensoring this identity derivedly with $C$ preserves the splitting
as a statement about $C$-complexes.  The total homology in degree
$i+h$ therefore contains $\Tor_i^R(C,L)$ as a direct summand.  If
$i>l=d-h$, that total degree exceeds $d$, proving the bound.  This
argument also shows directly that all the displayed lengths are
finite.
\end{proof}

\subsection{Finite approximations preserving the actions}

We need finite-length $B$-modules whose normalized Tor lengths approach
those of $U$, while retaining the single annihilator
$\m_B^a$.  No finite generation or finite presentation of $U$ is
assumed.  The approximation will take place among modules carrying
both its tower action and its $B$-action.  In this subsection,
$\lambda$ denotes $\lambda_\infty$ on underlying tower modules,
as permitted by Theorem~\ref{len:input}; the approximating modules
need not carry a $C$-action.

Let $A=A_D$, let $A_j$ be its flat regular tower from
Lemma~\ref{len:tower}, and let $\Lambda=\bigcup_j A_j$.
Thus
\[
 s(j):=\rank_A A_j=p^{jh}.
\]
Set
\begin{equation}\label{pos:approximation-algebras}
 S_j=A_j\otimes_A(D'/I_\sigma),\qquad
 S=\Lambda\otimes_A(D'/I_\sigma).
\end{equation}
The commuting $\Lambda$- and $D'/I_\sigma$-actions on $U$ combine
into an $S$-action.  Working with $S$-modules retains both of these
actions.  Descent to $S_j$ retains the existing $B$-action by
restriction along $B\to D'/I_\sigma$ induced by $\iota_\sigma$;
no new $B$-action is chosen at a finite stage.

\begin{lemma}\label{pos:approximation}
There are finite-length $B$-modules $M_n$ and positive integers
$s_n\to\infty$, with $\m_B^aM_n=0$, such that
\begin{align}
 \frac{\length_B(M_n)}{s_n}&\longrightarrow\lambda(U),
       \label{pos:approximation-length}\\
 \frac{\length_B\Tor_i^B(M_n,L)}{s_n}
       &\longrightarrow\tau_i(L)
       \quad(L=k,E',\ i\ge0).
       \label{pos:approximation-tor}
\end{align}
Writing $\beta_0(M_n)=\dim_k(M_n/\m_BM_n)$, we also have
\begin{equation}\label{pos:generator-bound}
 \frac{\beta_0(M_n)}{s_n}\longrightarrow\nu_\sigma,
 \qquad
 \nu_\sigma\ge
 \frac{\lambda(U)}{\length_B(B/\m_B^a)}>0.
\end{equation}
\end{lemma}

\begin{proof}
The algebra $S$ is finitely presented as a $\Lambda$-module and
has finite normalized length.  Every $S_j$ is finite over the
twisted image of $B$: $D'$ is finite over the twisted image of
$R/P$, and $A_j$ is finite free over $A$.  Equation
\eqref{pos:uniform-annihilator} annihilates $S_j$ and $S$.
Moreover, $S_j$ is Artinian local with residue field $k$.  To see
locality, the tower parameter roots are nilpotent over the
Artinian quotient $D'/I_\sigma$, and reduction leaves the common
residue field $k$.  Thus finite $S_j$-modules have finite lengths
over both $B$ and $A_j$, and these lengths agree: a composition
series over $S_j$ has factors $k$, of length one for either ring.
This argument is unchanged if a residue-field identification is
composed with an automorphism.

Given $\epsilon>0$, choose a finitely generated $S$-submodule
$V\subset U$ with
$\lambda(U/V)<\epsilon$.  This follows from the supremum property
of normalized length: the $S$-span of a finitely generated
$\Lambda$-submodule is still finitely generated over $\Lambda$.
To obtain a finite presentation as well, choose a surjection $S^r\to V$
with kernel $K$.  Its kernel has finite length, bounded by
$r\lambda(S)$.  Choose a finitely generated $S$-submodule $K_0\subset K$ such
that $\lambda(K/K_0)<\epsilon$.  Then
\[
 P_\epsilon=S^r/K_0\longrightarrow U
\]
has kernel and cokernel of length less than $\epsilon$.
It is finitely presented over $S$ and over $\Lambda$, and is
still annihilated by $\m_B^a$.

Fix $B$-resolutions of $k$ and $E'$ with finite free terms, and write
$r_i(L)$ for the rank in degree $i$, with $r_{-1}(L)=0$.
If $K_\epsilon$ and $W_\epsilon$ are the kernel and cokernel of
$P_\epsilon\to U$, the two long exact homology sequences give
\begin{align}
 &\left|\lambda\Tor_i^B(P_\epsilon,L)
                -\lambda\Tor_i^B(U,L)\right| \notag\\
 &\quad\le
   \bigl(r_i(L)+r_{i-1}(L)\bigr)\lambda(K_\epsilon)
   +\bigl(r_{i+1}(L)+r_i(L)\bigr)\lambda(W_\epsilon).
       \label{pos:tor-approximation-error}
\end{align}
For example, each homology of the complex obtained by tensoring
the resolution with $K_\epsilon$ has length at most the rank of
the corresponding term times $\lambda(K_\epsilon)$; the same
holds for $W_\epsilon$.  All these quantities are finite.  Also
$|\lambda(P_\epsilon)-\lambda(U)|<2\epsilon$.

A finite $S$-presentation of $P_\epsilon$ descends to an
$S_j$-module $P_j$ for some $j$, with
$P_\epsilon=\Lambda\otimes_{A_j}P_j$.  Flatness of
$\Lambda/A_j$ gives, for every $i$,
\[
 \Tor_i^B(P_\epsilon,L)
       =\Lambda\otimes_{A_j}\Tor_i^B(P_j,L).
\]
The finite-stage homology on the right is finite over $S_j$.
The normalization formula in Theorem~\ref{len:input}, and the
length comparison above, therefore give exactly
\[
 \lambda\Tor_i^B(P_\epsilon,L)
    =\frac{\length_{A_j}\Tor_i^B(P_j,L)}{s(j)}
    =\frac{\length_B\Tor_i^B(P_j,L)}{s(j)}.
\]
The same formula holds for $P_\epsilon$ itself.

At step $n$, choose $\epsilon$ so that all errors in
Equation~\eqref{pos:tor-approximation-error}, for $i\le n$ and
$L=k,E'$, and the length error are less than $1/n$.
Descend this presentation and enlarge its stage if necessary so
that $j\ge n$.  Further flat tower extension leaves every
normalized value unchanged.  Take $M_n=P_j$ and $s_n=s(j)$.
This proves Equations~\eqref{pos:approximation-length} and
\eqref{pos:approximation-tor}, including the common annihilator
and $s_n\to\infty$.

The case $i=0,L=k$ gives the limit in
Equation~\eqref{pos:generator-bound}.  A minimal generating set and the
annihilator give a surjection
$(B/\m_B^a)^{\beta_0(M_n)}\to M_n$.  Consequently
\[
 \length_B(M_n)
    \le\beta_0(M_n)\length_B(B/\m_B^a).
\]
Divide by $s_n$ and take limits to obtain the stated positive
lower bound.
\end{proof}

\subsection{Bounded complexes and their ordinary evaluation}

The approximations provide two kinds of control.  In minimal
resolutions, the Tor bound $\tau_i(k)=0$ for $i>l$ makes the
corresponding ranks negligible.  The common annihilator then makes
the truncated sequences
supported after quotienting by negligible ranks.  That same annihilator
has already given the positive normalized generator bound
\eqref{pos:generator-bound}, which will make the second evaluation
strictly positive.

Choose minimal free $B$-resolutions $F_n^\infty$ of $M_n$, indexed
homologically in nonnegative degrees.  Let $F_n$ be their brutal
truncations to degrees $0,\ldots,l$.  Minimality and
Equation~\eqref{pos:approximation-tor} give, in each fixed degree,
\begin{equation}\label{pos:rank-estimates}
 \rank_B F_{n,i}^\infty=O(s_n),\qquad
 \rank_B F_{n,i}^\infty=o(s_n)\quad(i>l),
\end{equation}
because these ranks equal $\length_B\Tor_i^B(M_n,k)$.

Let $z_1,\ldots,z_l$ be regular parameters of $A_E$.  For
$\gamma\in G_E$, extend $F_n$ through
\[
 B\longrightarrow E\longrightarrow E'
                  \xrightarrow{\gamma}E',
\]
and denote the resulting complex by $F_{n,\gamma}'$.
Use the normalization $s_n$ in the asymptotic category over $E'$.
These numbers come from the $D$-side approximation. Section~\ref{sec:asymptotic}
allows any positive integer sequence $s_n$, so no identification or
compatibility between the parameter-root towers of $A_D$ and $A_E$
is needed.

\begin{lemma}\label{pos:supported-truncation}
The sequence $(F_{n,\gamma}')_n$ is supported at the closed point
in that category.  Its ordinary Euler evaluation is $c_\sigma$
for every $\gamma\in G_E$.
\end{lemma}

\begin{proof}
Every element $b\in\m_B^a$ acts nullhomotopically on
$F_n^\infty$, because it annihilates $M_n$.  Restrict a chosen
homotopy to the brutal truncation.  Its only remaining error is
in degree $l$, where it is
\begin{equation}\label{pos:truncation-error}
 F_{n,l}\xrightarrow{h_l}F_{n,l+1}^\infty
                      \xrightarrow{d_{l+1}}F_{n,l}.
\end{equation}
This is a factorization by chain maps through
$F_{n,l+1}^\infty[l]$: the map induced by $d_{l+1}$ is a chain map
because $d_ld_{l+1}=0$, and the map induced by $h_l$ is a chain map
because the truncation has no degree $l+1$ term.  The middle module has rank
$o(s_n)$ by Equation~\eqref{pos:rank-estimates}, and is therefore killed
in the asymptotic quotient.  These identities persist after
extension to $E'$, without a flatness assumption over $B$.

The ideal $\m_BE'$ for each indicated structural map is primary
to $\m_{E'}$.  Choose $v$ so large that every $z_j^v$ belongs to
the extension of $\m_B^a$ for all $\gamma$.  This is one fixed
choice, since the group is finite.  Expressing $z_j^v$ as a
finite $E'$-linear combination of elements of $\m_B^a$ proves
that it acts by zero on the sequence in the quotient.  The
same combinations give genuine nullhomotopies on the full
extended resolutions.  This proves support.

Apply the ordinary version of
Equation~\eqref{as:koszul-polynomial}.  In its right-hand side, one may
replace the truncated extended resolutions by the full ones in
each fixed homological degree.  Indeed, the quotient of the
full resolution by its brutal truncation has only terms of
degrees at least $l+1$, all of rank $o(s_n)$ in each fixed degree.
After tensoring that tail with $K(z^v;E')$, taking Koszul
homology first bounds each fixed-degree homology by finite sums
of
\[
 \rank_B F_{n,i}^\infty\,
     \length_{E'} H_q(z^v;E')=o(s_n).
\]
The Koszul homology lengths are fixed and finite, and only
finitely many terms contribute to each total degree.  The long
exact homology sequence thus proves the replacement assertion.

On each full extended resolution the powers $z_j^v$ act
nullhomotopically, as established above.  Koszul tensoring
therefore gives the direct sum of shifts with multiplicities
$\binom lj$.  Its underlying homology has ordinary lengths
$\length_B\Tor_i^B(M_n,E')$, unchanged by conjugating the
$E'$-action: applying $\gamma$ entrywise is a semilinear
isomorphism.  Equations~\eqref{pos:approximation-tor} and
\eqref{pos:tor-identification}, followed by coefficient
comparison in Equation~\eqref{as:koszul-polynomial}, now identify the
ordinary homology-length coefficients with
\[
 \tau_i(E')=\lambda\Tor_i^R(C,E').
\]
Their alternating sum is exactly $c_\sigma$.
\end{proof}

\subsection{A positive normalized evaluation}

The ordinary evaluation has recovered $c_\sigma$ on every conjugate.
We now evaluate the same supported object using $C_E$.  The parameter
Koszul identities will leave only its nonnegative degree-zero
coefficient, and minimality of the resolutions will bound that
coefficient below using \eqref{pos:generator-bound}.

Corollary~\ref{len:finite-length-base-change} gives
\begin{equation}\label{pos:residue-positive}
 0<\eta_E:=\lambda_E(C_E/\m_{E'}C_E)<\infty.
\end{equation}
We record a quantitative bound for use below.  Set
$q_E=(z_1,\ldots,z_l)E'$,
$r_E=\rank_{A_E}E'$, and
$t_E=\length_{E'}(E'/q_E)$.  The length input gives
$\lambda_E(C_E/q_EC_E)=r_E>0$.
A composition series of $E'/q_E$, tensored right-exactly with
$C_E$, shows
\begin{equation}\label{pos:residue-bound}
 r_E\le t_E\eta_E.
\end{equation}
At each step the image of the tensor of the submodule has
length at most that tensor's length; injectivity after tensoring
is not required.  Finiteness of $\eta_E$ follows from the
surjection $C_E/q_EC_E\to C_E/\m_{E'}C_E$.

\begin{lemma}\label{pos:positive-evaluation}
The $\lambda_E$ Euler evaluation of $(F_{n,\gamma}')_n$ is a
number $b_{\sigma,\gamma}$ satisfying
\begin{equation}\label{pos:strict-bound}
 b_{\sigma,\gamma}
    \ge\nu_\sigma\eta_E
    \ge
       \frac{\lambda_D(U)}{\length_B(B/\m_B^a)}
       \frac{r_E}{t_E}
    >0.
\end{equation}
\end{lemma}

\begin{proof}
Let $b_i$ be the nonnegative homology-length coefficients of
this evaluation in Equation~\eqref{as:koszul-polynomial}.  Consider the
finite double complex
\[
 K(z^v;E')\otimes_{E'}F_{n,\gamma}'
                         \otimes_{E'}C_E.
\]
Taking Koszul homology first, Theorem~\ref{len:input} makes all
positive Koszul rows have length zero.  The remaining row has
complex degrees $0,\ldots,l$.  Thus the normalized homology
lengths of this double complex vanish in total degrees greater
than $l$, as well as in negative degrees.  Equation
\eqref{as:koszul-polynomial} says that their polynomial is
$(1+t)^l\sum_i b_it^i$.  All $b_i$ are nonnegative.  Boundedness
and comparison of the least degree rule out negative-degree
coefficients.  Any $b_i>0$ with $i>0$ would then give a positive
coefficient in degree $i+l>l$, a contradiction.  Consequently
only $b_0$ can be nonzero, and the Euler evaluation equals $b_0$.

In degree zero, the displayed double complex has homology
\[
 C_E^{\beta_0(M_n)}\big/
 \left(\operatorname{im}(d_1)+(z_1^v,\ldots,z_l^v)
                                    C_E^{\beta_0(M_n)}\right).
\]
Since $l>0$, $d_1$ is the degree-one presentation differential
of $M_n$.  Minimality over $B$ puts all its entries in $\m_B$,
whose image under any twist is contained in $\m_{E'}$.
The base parameter powers are also in $\m_{E'}$.  There is
therefore a surjection from this degree-zero homology onto
\[
 (C_E/\m_{E'}C_E)^{\beta_0(M_n)}.
\]
Its normalized length is at least
$\beta_0(M_n)\eta_E/s_n$.  The degree-zero coefficient in
Equation~\eqref{as:koszul-polynomial} is $b_0$, so taking the ultralimit
and using the ordinary limit in Equation~\eqref{pos:generator-bound}
gives $b_0\ge\nu_\sigma\eta_E$.  Equations
\eqref{pos:generator-bound} and \eqref{pos:residue-bound} give
the explicit remaining inequalities in
Equation~\eqref{pos:strict-bound}.  Define $b_{\sigma,\gamma}=b_0$.
\end{proof}

\begin{proof}[Proof of Theorem~\ref{thm:main}]
By Lemma~\ref{pos:reduction}, use the domain pair fixed above.
For each $\sigma\in G_D$, carry out the construction with its
specified $R$-action.  Apply Theorem~\ref{as:norm} to the supported
sequence of
Lemma~\ref{pos:supported-truncation}, over $A_E\subset E'$.
It equates the sums of the ordinary and normalized evaluations
of all $G_E$-conjugates.  The former evaluation is $c_\sigma$
for every conjugate, whereas the latter is
$b_{\sigma,\gamma}$.  Hence
\begin{equation}\label{pos:second-average}
 |G_E|c_\sigma
       =\sum_{\gamma\in G_E}b_{\sigma,\gamma}>0
\end{equation}
by Lemma~\ref{pos:positive-evaluation}.
Finally, Lemma~\ref{pos:first-average} gives
\[
 |G_D|a_Da_E\,\chi^R(D,E)
    =\frac{1}{|G_E|}
       \sum_{\sigma\in G_D}\sum_{\gamma\in G_E}
                   b_{\sigma,\gamma}>0.
\]
Both ranks and both group orders are positive.  Thus
$\chi^R(D,E)>0$, and the reduction proves the theorem.
Neither norm comparison requires a Galois action on $C_D$ or
$C_E$, nor invariance of either normalized length under a twist.
\end{proof}

\begin{thebibliography}{99}
\raggedright

\bibitem{BalmerSchlichting}
P.~Balmer and M.~Schlichting,
Idempotent completion of triangulated categories,
\emph{J. Algebra} \textbf{236} (2001), 819--834.
\href{https://doi.org/10.1006/jabr.2000.8529}{doi:10.1006/jabr.2000.8529}.

\bibitem{Gabber}
P.~Berthelot,
Alt\'erations de vari\'et\'es alg\'ebriques [d'apr\`es A.~J.~de~Jong],
\emph{Ast\'erisque} \textbf{241} (1997), Exp.~815, 273--311;
Section~6.1, Theorem~6.1.1 (Gabber).
\url{https://www.numdam.org/item/SB_1995-1996__38__273_0.pdf}.

\bibitem{BHM}
B.~Bhatt, M.~Hochster, and L.~Ma,
Lim Cohen--Macaulay sequences of modules,
\emph{J. Reine Angew. Math.} \textbf{835} (2026), 111--169.
Numbered references follow
\href{https://arxiv.org/abs/2410.18372v2}{arXiv:2410.18372v2},
August~25, 2026.

\bibitem{BhattScholze}
B.~Bhatt and P.~Scholze,
Prisms and prismatic cohomology,
\emph{Ann. of Math. (2)} \textbf{196} (2022), no.~3, 1135--1275.
\href{https://doi.org/10.4007/annals.2022.196.3.5}{doi:10.4007/annals.2022.196.3.5}.
Numbered references follow
\href{https://arxiv.org/html/1905.08229v4}{arXiv:1905.08229v4},
January~12, 2022.

\bibitem{BGT}
A.~J.~Blumberg, D.~Gepner, and G.~Tabuada,
A universal characterization of higher algebraic $K$-theory,
\emph{Geom. Topol.} \textbf{17} (2013), 733--838.
\href{https://doi.org/10.2140/gt.2013.17.733}{doi:10.2140/gt.2013.17.733}.
Numbered references follow
\href{https://arxiv.org/abs/1001.2282v4}{arXiv:1001.2282v4},
February~5, 2013.

\bibitem{CLMST}
H.~Cai, S.~Lee, L.~Ma, K.~Schwede, and K.~Tucker,
Perfectoid signature, perfectoid Hilbert--Kunz multiplicity, and an
application to local fundamental groups,
\emph{Invent. Math.} \textbf{244} (2026), 357--454.
\href{https://doi.org/10.1007/s00222-025-01399-w}{doi:10.1007/s00222-025-01399-w}.
Numbered references follow
\href{https://arxiv.org/html/2209.04046v4}{arXiv:2209.04046v4},
July~3, 2025.

\bibitem{Drinfeld}
V.~Drinfeld,
DG quotients of DG categories,
\emph{J. Algebra} \textbf{272} (2004), no.~2, 643--691.
\href{https://doi.org/10.1016/j.jalgebra.2003.05.001}{doi:10.1016/j.jalgebra.2003.05.001}.
Numbered references follow
\href{https://arxiv.org/abs/math/0210114v7}{arXiv:math/0210114v7},
February~17, 2008.

\bibitem{DuttaFiniteSupport}
S.~P.~Dutta,
A special case of positivity (II),
\emph{Proc. Amer. Math. Soc.} \textbf{133} (2005), no.~7, 1891--1896.
\href{https://doi.org/10.1090/S0002-9939-05-07929-3}{doi:10.1090/S0002-9939-05-07929-3}.

\bibitem{Faltings}
G.~Faltings,
Almost \'etale extensions,
\emph{Ast\'erisque} \textbf{279} (2002), 185--270.
\href{https://numdam.org/item/AST_2002__279__185_0/}{Numdam: AST\_2002\_\_279\_\_185\_0}.

\bibitem{GabberRamero}
O.~Gabber and L.~Ramero,
\emph{Foundations for almost ring theory---Release 7.5},
\href{https://arxiv.org/abs/math/0409584v13}{arXiv:math/0409584v13},
September~29, 2018.

\bibitem{GilletSoule}
H.~Gillet and C.~Soul\'e,
Intersection theory using Adams operations,
\emph{Invent. Math.} \textbf{90} (1987), no.~2, 243--277.
\href{https://doi.org/10.1007/BF01388705}{doi:10.1007/BF01388705}.

\bibitem{Goodwillie}
T.~G.~Goodwillie,
Relative algebraic $K$-theory and cyclic homology,
\emph{Ann. of Math. (2)} \textbf{124} (1986), no.~2, 347--402.
\href{https://doi.org/10.2307/1971283}{doi:10.2307/1971283}.

\bibitem{IshiroShimomoto}
S.~Ishiro and K.~Shimomoto,
$\delta$-rings, perfectoid towers, and lim Cohen--Macaulay sequences,
preprint,
\href{https://arxiv.org/abs/2509.06527v3}{arXiv:2509.06527v3},
May~31, 2026.

\bibitem{Jannsen}
U.~Jannsen,
Resolution of singularities for embedded curves,
Appendix~A to S.~Saito and K.~Sato,
A finiteness theorem for zero-cycles over $p$-adic fields,
\emph{Ann. of Math. (2)} \textbf{172} (2010), no.~3, 1593--1639,
Proposition~A.4(1).
\url{https://annals.math.princeton.edu/wp-content/uploads/annals-v172-n3-p03-p.pdf}.

\bibitem{KCSotoLevins}
N.~KC and A.~J.~Soto Levins,
On liftings of modules of finite projective dimension,
\emph{Int. Math. Res. Not. IMRN} \textbf{2024}, no.~24, 14729--14736.
\href{https://doi.org/10.1093/imrn/rnae257}{doi:10.1093/imrn/rnae257}.
Numbered references follow
\href{https://arxiv.org/abs/2404.17572v2}{arXiv:2404.17572v2},
November~5, 2024.

\bibitem{KuranoRoberts}
K.~Kurano and P.~C.~Roberts,
The positivity of intersection multiplicities and symbolic powers of
prime ideals,
\emph{Compositio Math.} \textbf{122} (2000), no.~2, 165--182.
\href{https://doi.org/10.1023/A:1001741230307}{doi:10.1023/A:1001741230307}.

\bibitem{CohenGabber}
K.~Kurano and K.~Shimomoto,
An elementary proof of Cohen--Gabber theorem in the equal characteristic
$p>0$ case,
\emph{Tohoku Math. J. (2)} \textbf{70} (2018), no.~3, 377--389,
Theorem~1.1.
\href{https://arxiv.org/abs/1510.03573}{arXiv:1510.03573}.

\bibitem{LT}
M.~Land and G.~Tamme,
On the $K$-theory of pullbacks,
\emph{Ann. of Math. (2)} \textbf{190} (2019), no.~3, 877--930.
\href{https://doi.org/10.4007/annals.2019.190.3.4}{doi:10.4007/annals.2019.190.3.4}.
Numbered references follow
\href{https://arxiv.org/html/1808.05559v3}{arXiv:1808.05559v3},
November~8, 2019.

\bibitem{Ma}
L.~Ma,
Lech's conjecture in dimension three,
\emph{Adv. Math.} \textbf{322} (2017), 940--970.
\href{https://doi.org/10.1016/j.aim.2017.10.032}{doi:10.1016/j.aim.2017.10.032}.
Numbered references follow
\href{https://arxiv.org/html/1609.00095v3}{arXiv:1609.00095v3},
February~13, 2018.

\bibitem{Monsky}
P.~Monsky,
The Hilbert--Kunz function,
\emph{Math. Ann.} \textbf{263} (1983), no.~1, 43--49.
\href{https://doi.org/10.1007/BF01457082}{doi:10.1007/BF01457082}.

\bibitem{OpenAILech2026}
OpenAI,
Lech's multiplicity conjecture,
OpenAI Math Release preprint
\href{https://github.com/openai/math/blob/main/preprints/Lechs-multiplicity-conjecture-September-23-2026/paper.pdf}{OAI:Lechs-multiplicity-conjecture-September-23-2026},
2026.

\bibitem{Roberts}
P.~Roberts,
The vanishing of intersection multiplicities of perfect complexes,
\emph{Bull. Amer. Math. Soc. (N.S.)} \textbf{13} (1985), no.~2, 127--130.
\href{https://doi.org/10.1090/S0273-0979-1985-15394-7}{doi:10.1090/S0273-0979-1985-15394-7}.

\bibitem{RobertsRennes}
P.~Roberts,
Intersection multiplicities in commutative algebra,
\emph{Publications de l'Institut de recherche math\'ematiques de Rennes}
(1985), fasc.~4, 25--39.
\url{https://www.numdam.org/article/PSMIR_1985___4_25_0.pdf}.

\bibitem{RobertsIntersectionTheorems}
P.~Roberts,
Intersection theorems,
in \emph{Commutative Algebra} (Berkeley, CA, 1987),
M.~Hochster, C.~Huneke, and J.~D.~Sally (eds.),
Mathematical Sciences Research Institute Publications \textbf{15},
Springer, New York, 1989, 417--436.
\href{https://doi.org/10.1007/978-1-4612-3660-3_23}{doi:10.1007/978-1-4612-3660-3\_23}.

\bibitem{Serre}
J.-P.~Serre,
\emph{Local Algebra},
translated by C.~W.~Chin,
Springer Monographs in Mathematics,
Springer, Berlin, 2000, Chapter~V.
\href{https://doi.org/10.1007/978-3-662-04203-8}{doi:10.1007/978-3-662-04203-8}.

\bibitem{Skalit}
C.~Skalit,
Positivity of intersection multiplicity over a two-dimensional base,
\emph{J. Pure Appl. Algebra} \textbf{223} (2019), no.~4, 1801--1816.
\href{https://doi.org/10.1016/j.jpaa.2018.07.008}{doi:10.1016/j.jpaa.2018.07.008}.
Numbered references follow
\href{https://arxiv.org/abs/1510.05146v2}{arXiv:1510.05146v2},
July~31, 2018.

\bibitem{SkalitUnramified}
C.~Skalit,
Intersection multiplicity of Serre in the unramified case,
\emph{J. Algebra} \textbf{428} (2015), 91--121.
\href{https://doi.org/10.1016/j.jalgebra.2014.12.040}{doi:10.1016/j.jalgebra.2014.12.040}.
Numbered references follow
\href{https://arxiv.org/abs/1409.3616v2}{arXiv:1409.3616v2},
December~10, 2014.

\bibitem{Stacks}
The Stacks Project Authors,
\emph{The Stacks Project},
\href{https://stacks.math.columbia.edu/tag/0815}{Tag 0815},
\href{https://stacks.math.columbia.edu/tag/03QH}{Tag 03QH},
\href{https://stacks.math.columbia.edu/tag/0BRI}{Tag 0BRI}, and
\href{https://stacks.math.columbia.edu/tag/0BRJ}{Tag 0BRJ},
accessed September~23, 2026.

\bibitem{ThomasonTrobaugh}
R.~W.~Thomason and T.~Trobaugh,
Higher algebraic $K$-theory of schemes and of derived categories,
in \emph{The Grothendieck Festschrift}, Vol.~III,
Progress in Mathematics \textbf{88}, Birkh\"auser Boston, 1990, 247--435.
\href{https://www.sas.rochester.edu/mth/sites/doug-ravenel/otherpapers/Thomason-Trobaugh.pdf}{Original-edition text}.

\end{thebibliography}

\end{document}
